\section{Symmetry examples, time-reversal parity, and the finite Heisenberg action}

\subsection{Symmetry examples}

\begin{lemma}[All phases under conjugation]
    \label{ex:all-phases-under-conjugation}
    \notready
    \discussion{5715}
    \uses{conjphs, prop:conjsym-strict-equiv, corconj, cor:sympl-to-orthogonal}
    Let $u_+=v_+=\Id_{r\ell}$ and
    $u_-=v_-=\operatorname{diag}(-1,1,\ldots,1)$. If at least one of $r$ and
    $\ell$ is odd, $(u_+,v_+)$ and $(u_+,v_-)$ realize the two classes of
    Theorem~\ref{conjphs}. If both are even, the four pairs $(u_\pm,v_\pm)$
    realize the four Case I classes. Multiplying by the matrix $R$ of
    Lemma~\ref{cor:sympl-to-orthogonal} realizes the four Case II classes.
    % Source lines: paper_v2.tex 1845--1876.
\end{lemma}

\begin{lemma}[Admissible phases under conjugation]
    \label{ex:mpu_admissible_all_conjugation_phases}
    \notready
    \discussion{5715}
    % **Scope restriction (realization by source-derived standard forms):** each
    % displayed standard-form tensor carries reduced full-support source data.
    % Documented in docs/paper-gaps/mpu_canonical_form_full_support.tex.
    \uses{thm:mpu_admissible_conjugation_odd_rank, thm:mpu_admissible_conjugation_even_rank, prop:mpu_admissible_conjugation_standard_forms, cor:mpu_admissible_symplectic_to_orthogonal, def:mpu_reduced_full_support_source_datum, def:mpu_admissible_equivalence_datum}
    Assume the standard-form tensors displayed below carry reduced full-support
    source data and that the comparison paths used to classify them carry
    admissible path data of blocking length one. Let
    \begin{align}
        u_+ & =v_+=\Id_{r\ell},                        &
        u_- & =v_-=\operatorname{diag}(-1,1,\ldots,1).
    \end{align}
    If at least one of $r$ and $\ell$ is odd, the pairs $(u_+,v_+)$ and
    $(u_+,v_-)$ represent the two classes of Lemma~\ref{thm:mpu_admissible_conjugation_odd_rank}. If $r$
    and $\ell$ are both even, the four pairs $(u_\pm,v_\pm)$ represent the
    four Case I classes of Lemma~\ref{thm:mpu_admissible_conjugation_even_rank}. With
    \begin{align}
        u'_\pm & =(\Id_{r\ell/4}\otimes R)u_\pm, &
        v'_\pm & =v_\pm(\Id_{r\ell/4}\otimes R),
    \end{align}
    the four pairs $(u'_\pm,v'_\pm)$ represent the four Case II classes.
    % Source lines: paper_v2.tex 1845--1876.
\end{lemma}

\begin{lemma}[Normalization obstruction in the printed CZX tensor]
    \label{thm:czx_printed_normalization_obstruction}
    % Principal declaration: MPOTensor.CZX.printed2017Tensor_not_isMPU.
    \lean{MPOTensor.CZX.printed2017Tensor, MPOTensor.CZX.printed2017Tensor_two_zero_row, MPOTensor.CZX.printed2017Tensor_not_isMPU}
    \leanok
    \discussion{7738}
    \uses{def:mpu}
    Let $u,d,\ell,r\in\{0,1\}$ and set
    $P^{ud}_{\ell r}=\delta_{u,1-d}\delta_{u,\ell}2^{-1/2}(-1)^{\ell r}$
    as in the CZX example of \cite{Cirac2017MPU} for the printed local tensor.
    In the ordinary periodic trace contraction, the two-site operator
    generated by $P$ has all-zero output row
    \begin{align}
        \operatorname{MPO}_2(P)_{00,b} & =\tfrac12\delta_{b,11}.
    \end{align}
    In particular, $P$ does not generate a matrix product unitary.
    The note \cite{gap:mpu_czx_tensor_normalization} details the normalization discrepancy and its comparison with the Z-decorated
    tensor.
    % Source: paper_v2.tex, example:czx, lines 1914--1930.
    % False-source note: docs/paper-gaps/mpu_czx_tensor_normalization.tex.
\end{lemma}
\begin{proof}
    \leanok
    \uses{lem:mpu_mul_adjoint}
    The copy deltas fix both virtual indices to zero. The physical bit
    flips then force the input to be $11$, and the two Hadamard entries
    contribute $1/2$. If the operator were unitary, its all-zero diagonal
    entry after multiplication by its adjoint would give
    \begin{align}
        1 & =\sum_b\left|\operatorname{MPO}_2(P)_{00,b}\right|^2
        =\left(\tfrac12\right)^2=\tfrac14,
    \end{align}
    a contradiction.
\end{proof}

\begin{remark}
    The normalized Hadamard in the following printed CZX formulas has
    the obstruction above. In contrast, in \cite{FrancoRubio2025MPUGauging} the blue tensors of
    the CZX tensor diagram
    are $\sqrt{2}H$, and this diagram also has a $Z$ gate on each
    output qubit. Let $A=\sqrt2 P$, let $A_2$ denote its two-site block,
    and let $B$ denote that displayed Z-decorated tensor. The output-leg
    relation is
    \begin{align}
        B & =(Z\otimes Z)A_2.
        \label{eq:mpu_czx_decorated_block}
    \end{align}
    The note \cite{gap:mpu_czx_tensor_normalization} gives the coordinate derivation of
    (\ref{eq:mpu_czx_decorated_block}).
\end{remark}

\begin{lemma}[CZX]
    \label{example:czx}
    \notready
    \discussion{5715}
    \uses{prop:sigma1-and-sigma2-equal, prop:dagger-mps-class-and-mpu-class-equal, thm:czx_printed_normalization_obstruction}
    The normalized tensor printed below is not an MPU, by
    Lemma~\ref{thm:czx_printed_normalization_obstruction}.
    The following is the source classification claim, not an assertion
    of unitarity for that printed tensor; \cite{gap:mpu_czx_tensor_normalization} discusses its normalization and comparison
    with the Z-decorated
    tensor.
    The CZX MPU of Chen, Liu and Wen~\cite{ChenLiuWen2011CZX} (see also
    \cite[lines~1906--1943]{Cirac2017MPU}) is defined by
    \begin{align}
        u & =CZ(\sigma_x\otimes\sigma_x), & v & =CZ,
    \end{align}
    where $CZ=\operatorname{diag}(1,1,1,-1)$. Let
    \begin{align}
        H & =\frac1{\sqrt2}\begin{pmatrix}1&1\\1&-1\end{pmatrix}
    \end{align}
    be the Hadamard matrix. The corresponding tensor is
    \begin{align}
        \mathcal U_{\ell r}^{ud}
         & =(\sigma_x)_{ud}\delta_{u,\ell}H_{\ell,r}.
    \end{align}
    It has $\sigma^{(1)}=\sigma^{(2)}=-1$.
    % Source lines: paper_v2.tex 1906--1943.
\end{lemma}

\begin{lemma}[Admissible CZX data]
    \label{ex:mpu_admissible_czx}
    \notready
    \discussion{5715}
    % **Scope restriction (blocked CZX data):** the one-, two-, and four-site
    % tensors used to define the two signs carry reduced full-support source data.
    % Documented in docs/paper-gaps/mpu_canonical_form_full_support.tex.
    \uses{prop:mpu_admissible_sigma_blocking_parity, prop:mpu_admissible_mps_time_reversal_class, def:mpu_reduced_full_support_source_datum, def:mpu_admissible_equivalence_datum, thm:czx_printed_normalization_obstruction}
    The normalized tensor printed below is not an MPU, by
    Lemma~\ref{thm:czx_printed_normalization_obstruction}.
    Under the conventions recorded in \cite{gap:mpu_czx_tensor_normalization}, the classification below evaluates the given tensor
    data.
    Suppose the CZX tensors $\mathcal U$, $\mathcal U_2$, and $\mathcal U_4$,
    together with the physical-adjoint comparison tensors
    $\mathcal U_2^\sharp$ and $\mathcal U_4^\sharp$, are simple and each carries
    its own reduced full-support source datum. The CZX MPU of
    Chen, Liu and Wen~\cite{ChenLiuWen2011CZX} (see also
    \cite[lines~1906--1943]{Cirac2017MPU}) is defined by
    \begin{align}
        u  & =CZ(\sigma_x\otimes\sigma_x),   & v & =CZ, \\
        CZ & =\operatorname{diag}(1,1,1,-1).
    \end{align}
    Its tensor is
    \begin{align}
        \mathcal U_{\ell r}^{ud}
          & =(\sigma_x)_{ud} \delta_{u,\ell}H_{\ell,r},           &
        H & =\frac1{\sqrt2}\begin{pmatrix}1&1\\1&-1\end{pmatrix}.
    \end{align}
    It has $\sigma^{(1)}=\sigma^{(2)}=-1$.
    % Source lines: paper_v2.tex 1906--1943.
\end{lemma}

\begin{lemma}[The other three time-reversal classes]
    \label{ex:dagger-all-but-czx}
    \notready
    \discussion{5715}
    \uses{example:czx, prop:sigma1-and-sigma2-equal}
    The standard forms
    \begin{align}
        u & =v=\Id,                                      \\
        u & =i\Id,                           & v & =\Id, \\
        u & =i\,CZ(\sigma_x\otimes\sigma_x), & v & =CZ
    \end{align}
    realize, respectively,
    \begin{align}
        (\sigma^{(1)},\sigma^{(2)})
         & =(1,1),\quad(-1,1),\quad(1,-1).
    \end{align}
    % Source lines: paper_v2.tex 1945--1962.
\end{lemma}

\begin{lemma}[Admissible remaining time-reversal classes]
    \label{ex:mpu_admissible_time_reversal_classes}
    \notready
    \discussion{5715}
    % **Scope restriction (blocked example data):** every standard form and
    % blocking used to define the signs carries reduced full-support source data.
    % Documented in docs/paper-gaps/mpu_canonical_form_full_support.tex.
    \uses{ex:mpu_admissible_czx, prop:mpu_admissible_sigma_blocking_parity, def:mpu_reduced_full_support_source_datum, def:mpu_admissible_equivalence_datum}
    For each displayed example, suppose the tensors $\mathcal U$,
    $\mathcal U_2$, and $\mathcal U_4$, together with the physical-adjoint
    comparison tensors $\mathcal U_2^\sharp$ and $\mathcal U_4^\sharp$, are
    simple and each carries its own reduced full-support source datum.
    The standard forms
    \begin{align}
        u & =v=\Id,                                      \\
        u & =i\Id,                           & v & =\Id, \\
        u & =i\,CZ(\sigma_x\otimes\sigma_x), & v & =CZ
    \end{align}
    realize, respectively,
    \begin{align}
        (\sigma^{(1)},\sigma^{(2)})
         & =(1,1),\quad(-1,1),\quad(1,-1).
    \end{align}
    % Source lines: paper_v2.tex 1945--1962.
\end{lemma}

\begin{lemma}[Source cuts and ranks of the cyclic shifts]
    \label{thm:mpu_shift_source_cuts_ranks}
    \lean{MPOTensor.sourceCutM₁_rightShiftTensor}
    \lean{MPOTensor.sourceCutM₂_rightShiftTensor}
    \lean{MPOTensor.rightRank_rightShiftTensor}
    \lean{MPOTensor.leftRank_rightShiftTensor}
    \lean{MPOTensor.sourceCutM₁_leftShiftTensor}
    \lean{MPOTensor.sourceCutM₂_leftShiftTensor}
    \lean{MPOTensor.rightRank_leftShiftTensor}
    \lean{MPOTensor.leftRank_leftShiftTensor}
    \leanok
    \uses{def:mpu_shift_tensors, def:mpu_source_matrix_one, def:mpu_source_matrix_two, def:mpu_right_rank, def:mpu_left_rank}
    Let $A$ be the right-shift tensor and let
    $z_{(\alpha,i)}=\delta_{\alpha i}$ be the vectorized identity. Its
    source cuts are
    \begin{align}
        M_1(A) & =\Id,            \\
        M_2(A) & =zz^{\mathsf T}.
    \end{align}
    For the left-shift tensor $A^\sharp$ they are
    \begin{align}
        M_1(A^\sharp) & =zz^{\mathsf T}, \\
        M_2(A^\sharp) & =\Id.
    \end{align}
    In every physical dimension, including $d=0$,
    \begin{align}
        r(A)           & =d^2, \\
        \ell(A^\sharp) & =d^2.
    \end{align}
    If $d>0$, the remaining two ranks are
    \begin{align}
        \ell(A)     & =1, \\
        r(A^\sharp) & =1.
    \end{align}
    These are the source-cut quantities from
    \cite[Section~III]{Cirac2017MPU}; no choice of source factors is involved.
\end{lemma}

\begin{proof}\leanok
    \uses{def:mpu_shift_tensors, def:mpu_source_matrix_one, def:mpu_source_matrix_two, def:mpu_right_rank, def:mpu_left_rank, thm:mpu_physical_adjoint_source_cuts, thm:mpu_physical_adjoint_source_ranks}
    The entry formula $A^{ij}_{\alpha\beta}
        =\delta_{i\alpha}\delta_{j\beta}$ gives $M_1(A)=\Id$ and
    $M_2(A)=zz^{\mathsf T}$. Hence the first rank is $d^2$. For $d>0$, the
    vector $z$ has a nonzero diagonal coordinate, so its outer product has
    rank one. Physical adjunction exchanges the two cuts and their ranks,
    which gives all four left-shift formulas without a second rank
    calculation.
\end{proof}

\begin{lemma}[Unitarity of the swap operator]
    \label{thm:mpu_swap_matrix_unitary}
    \lean{Matrix.swapMatrix_isUnitaryBetween}
    \leanok
    Let $d\in\N$. On $\C^d\otimes\C^d$, let $\mathbb S$ be the swap operator
    used in \cite[Section~VII]{Cirac2017MPU}, defined on the standard basis by
    $\mathbb S(e_k\otimes e_\ell)=e_\ell\otimes e_k$. With columns indexed by
    source coordinates $(k,\ell)$ and rows indexed by target coordinates
    $(i,j)$, its entries are
    \begin{align}
        \mathbb S_{(i,j),(k,\ell)} & =\delta_{i,\ell}\delta_{j,k}.
    \end{align}
    Then $\mathbb S$ is unitary between the two product-coordinate orders:
    \begin{align}
        \mathbb S^\dagger\mathbb S & =\Id, \\
        \mathbb S\mathbb S^\dagger & =\Id.
    \end{align}
    % Operator characterization: paper_v2.tex lines 1316--1320.
    % Section VII uses: lines 1999--2034, equations eq:SF_u1_u3,
    % eq:uv2_U2, and eq:uv2_U3; the unitary interpolation has endpoint
    % S at lines 2148--2150.
\end{lemma}

\begin{proof}\leanok
    The swap operator is self-adjoint and involutive, so
    $\mathbb S^\dagger=\mathbb S$ and $\mathbb S^2=\Id$. Substitution gives
    both identities.
\end{proof}

\begin{lemma}[The shift products $U_2$ and $U_3$ are strictly equivalent]
    \label{thm:mpu_shift_products_strictly_equivalent}
    \lean{MPOTensor.shiftExampleU₂_strictlyEquivalent_shiftExampleU₃,
        MPOTensor.shiftExampleU₂_isMPUCanonicalForm,
        MPOTensor.shiftExampleU₃_isMPUCanonicalForm,
        MPOTensor.IsMPU.conjugatePhysical, MPOTensor.IsMPU.bondConj}
    \leanok
    \uses{threeMPU, def:strictly-equivalent-tensors, thm:mpu_canonical_form_of_cfii,
        thm:mpu-continuous-unitary-swap-path}
    For $d\geq1$ the tensors $U_2=T^\dagger\otimes T$ and $U_3=T\otimes T^\dagger$ are in
    canonical form and strictly equivalent, with no bond directions adjoined.
    Conjugating the physical legs of a tensor by a unitary, or its bond space by a
    unitary, preserves the matrix product unitary property.
    % Source: paper_v2.tex, lines 1980--2001 and 2138--2151, where the source joins
    % the ancilla-dressed $\widetilde U_{2,3}$ by one physical path; here $U_2$ and
    % $U_3$ themselves are joined by a physical path followed by a bond path.
\end{lemma}

\begin{proof}\leanok
    \uses{thm:mpu_canonical_form_of_cfii, thm:mpu-continuous-unitary-swap-path}
    The letter of $U_3$ at the exchanged physical indices is the letter of $U_2$
    conjugated on the bond by the factor exchange $\mathbb S$. Let $\mathbb S(x)$,
    $x\in[0,1]$, be the continuous unitary path from the identity to $\mathbb S$.
    The path $x\mapsto\mathbb S(x)$ acting on the physical legs of $U_2$ ends at
    $U_2$ with its two spins exchanged; the path $x\mapsto\mathbb S(x)(\cdot)\mathbb S(x)^\dagger$
    on the bond space then ends at $U_3$. Every point generates matrix product
    unitaries, since $\mathbb S(x)^{\otimes N}U^{(N)}\mathbb S(x)^{\dagger\otimes N}$ is
    unitary and bond conjugation changes no operator. Canonical form of the
    endpoints: the shift is in canonical form II, so are its physical adjoint
    and the tensor products, hence they are in canonical form.
\end{proof}

\begin{definition}[Trace-normalized shift weight]
    \label{def:mpu_shift_trace_weight}
    \lean{MPOTensor.shiftPaperWeight,
        MPOTensor.shiftPaperWeightSquared,
        MPOTensor.shiftPaperWeightSquared_eq}
    \leanok
    \uses{def:mpu_source_weight, def:mpu_shift_tensors, threeMPU}
    Canonical form II requires the positive right fixed point $\rho$ to have
    trace one. If the virtual weight on the bond space of a single shift is
    chosen to be a scalar identity, this normalization forces
    \begin{align}
        \rho_d & =d^{-1}\Id_d
    \end{align}
    whenever $d>0$. The same formula is defined at $d=0$, although it then has
    trace zero. The tensors $U_2$ and $U_3$ have bond dimension $d^2$, so the
    corresponding scalar trace-one matrix is
    \begin{align}
        \rho_{d^2} & =(d^2)^{-1}\Id_{d^2}.
    \end{align}
    These are the scalar identity choices compatible with the source weight in
    \cite[equation labelled~\texttt{Y1Y1X1X1}]{Cirac2017MPU}. No transfer-map
    fixed-point equation is asserted here.

    The explicit source factorizations below instead use
    $\rho_{\mathrm{aux}}=\Id_d$ for each constituent shift and
    $\rho_{\mathrm{aux}}=\Id_{d^2}$ after tensoring. The paper gates computed
    from them carry reciprocal factors $d$ and $d^{-1}$ because these identity
    weights differ from $\rho_d$ and $\rho_{d^2}$ when $d>1$.
    % Source normalization: paper_v2.tex 269--280.
    % Source-factor weight: paper_v2.tex 487--495.
    % Shift and tensor-product dimensions: paper_v2.tex 1980--1994.
\end{definition}

\begin{lemma}[Positivity and normalization of the shift weight]
    \label{thm:mpu_shift_trace_weight_normalized}
    \lean{MPOTensor.shiftPaperWeight_posSemidef,
        MPOTensor.shiftPaperWeight_posDef,
        MPOTensor.shiftPaperWeight_trace,
        MPOTensor.shiftPaperWeightSquared_posSemidef,
        MPOTensor.shiftPaperWeightSquared_posDef,
        MPOTensor.shiftPaperWeightSquared_trace}
    \leanok
    \uses{def:mpu_shift_trace_weight}
    For every $d\in\mathbb N$, the matrix $\rho_d$ is positive semidefinite. If
    $d>0$, then $\rho_d$ is positive definite and
    \begin{align}
        \tr(\rho_d) & =1.
    \end{align}
    Under the same hypothesis, $\rho_{d^2}$ is positive definite with trace one;
    it is positive semidefinite for every $d$.
\end{lemma}

\begin{proof}\leanok
    \uses{def:mpu_shift_trace_weight}
    The scalar $d^{-1}$ is nonnegative, and it is strictly positive when
    $d>0$. Hence the claims follow from positivity of the identity matrix.
    For positive $d$,
    \begin{align}
        \tr(\rho_d) & =d^{-1}\tr(\Id_d)=d^{-1}d=1.
    \end{align}
    Replacing $d$ by $d^2$ proves the product-space statements.
\end{proof}

\begin{definition}[Trace-normalized source factors for the right shift]
    \label{def:mpu_right_shift_paper_source_factors}
    \lean{MPOTensor.rightShiftPaperSourceFactors}
    \leanok
    \uses{def:mpu_shift_trace_weight, def:threeMPU_supplied_source_factors}
    Assume $d>0$, and let $T$ be the right-shift tensor. Starting from the
    displayed source factors for $T$ with identity weight, rescale the first
    cut by
    \begin{align}
        X_1' & =\sqrt d\,X_1,
             & Y_1'           & =d^{-1/2}Y_1,
             & Z_1'           & =\sqrt d\,Z_1,
    \end{align}
    and leave $X_2,Y_2,Z_2$ unchanged. Then these are source factors
    for the trace-normalized scalar weight $\rho_d=d^{-1}\Id_d$ prescribed by
    CFII. In particular,
    \begin{align}
        M_i & =X_i'Y_i',
            & (X_1')^\dagger(\Id_d\otimes\rho_d)X_1' & =\Id,
            & (X_2')^\dagger X_2'                    & =\Id,
            & Y_i'Z_i'                               & =\Id.
    \end{align}
    This gives the source-factor metric normalization of
    \cite[equations~\texttt{Erightleft}, \texttt{X1X2b}, and
        \texttt{YZ=1}]{Cirac2017MPU} for the right shift.
\end{definition}

\begin{lemma}[The paper source tensor $u$ for the right shift]
    \label{thm:mpu_right_shift_paper_source_u}
    \lean{MPOTensor.sourceU_rightShiftPaperSourceFactors_apply,
        MPOTensor.sourceU_rightShiftPaperSourceFactors_gram,
        MPOTensor.sourceU_rightShiftPaperSourceFactors_isIsometry}
    \leanok
    \uses{def:mpu_right_shift_paper_source_factors, def:mpu_source_uv,
        thm:mpu_shift_source_cuts_ranks}
    Assume $d>0$. Identify the unique left source coordinate with $0$ and the
    right source coordinate with $(a,b)$. For the trace-normalized source
    factors of the right shift,
    \begin{align}
        u_{(0,(a,b)),(i_1,i_2)}
         & =\delta_{a,i_1}\delta_{b,i_2}.
    \end{align}
    Hence
    \begin{align}
        \sum_{l,r}u_{(l,r),q}\overline{u_{(l,r),p}}
         & =\delta_{p,q},
         & u^\dagger u    & =\Id.
    \end{align}
    This fixes the order and normalization of the two triangles in
    \cite[equations~\texttt{uu} and~\texttt{uUnitary}]{Cirac2017MPU} for this
    basic MPU. It is the explicit normalization and orientation regression for
    Lemma~\ref{thm:mpu_source_u_retained_interior} and
    Lemma~\ref{lem:mpu_admissible_source_u_isometry}.
\end{lemma}

\begin{proof}\leanok
    \uses{def:mpu_right_shift_paper_source_factors, def:mpu_source_uv}
    For the identity-weight factors, direct substitution in
    $u=Y_2\mathbin{-}Y_1$ gives
    \begin{align}
        u^{(0)}_{(0,(a,b)),(i_1,i_2)}
         & =d^{1/2}\delta_{a,i_1}\delta_{b,i_2}.
    \end{align}
    The trace-normalized first cut replaces $Y_1$ by $d^{-1/2}Y_1$ and leaves
    $Y_2$ unchanged. Therefore
    \begin{align}
        u_{(0,(a,b)),(i_1,i_2)}
         & =d^{-1/2}u^{(0)}_{(0,(a,b)),(i_1,i_2)}
        =\delta_{a,i_1}\delta_{b,i_2}.
    \end{align}
    Summing the product of two such entries over $(a,b)$ gives
    $\delta_{p,q}$, which is the entrywise identity $u^\dagger u=\Id$.
\end{proof}

\begin{lemma}[Local physical-species swap for the shift MPUs]
    \label{thm:threeMPU_local_physical_swap}
    \lean{MPOTensor.shiftExampleU₃_physicalSwap_eq_swapMatrix_mul_shiftExampleU₂_mul_swapMatrix}
    \leanok
    \uses{threeMPU, def:mpu_shift_tensors, def:mpu_independent_tensor_product}
    Write $[U_a]^{(i_1,i_2),(j_1,j_2)}$ for the virtual matrix of the
    one-site tensor $U_a$ with physical output coordinate $(i_1,i_2)$ and
    input coordinate $(j_1,j_2)$. Let $\mathbb S$ be the SWAP matrix on the
    two virtual factors, defined by
    \begin{align}
        \mathbb S\lvert\alpha_1,\alpha_2\rangle
         & =\lvert\alpha_2,\alpha_1\rangle.
    \end{align}
    Then
    \begin{align}
        [U_3]^{(i_2,i_1),(j_2,j_1)}
         & =\mathbb S[U_2]^{(i_1,i_2),(j_1,j_2)}\mathbb S.
    \end{align}
    Thus exchanging the two physical species exchanges the local tensors
    $U_2$ and $U_3$, with the same exchange on their virtual factors.
    % Local refinement of the finite-chain statements at paper_v2.tex
    % 1999--2001 and 2240--2245.
\end{lemma}

\begin{proof}\leanok
    \uses{threeMPU, def:mpu_shift_tensors, def:mpu_independent_tensor_product}
    In virtual coordinates $(\alpha_1,\alpha_2)$ and
    $(\beta_1,\beta_2)$, multiplication by $\mathbb S$ on both sides selects
    the entry with row $(\alpha_2,\alpha_1)$ and column
    $(\beta_2,\beta_1)$. Consequently,
    \begin{align}
        [U_3]^{(i_2,i_1),(j_2,j_1)}_{(\alpha_1,\alpha_2),(\beta_1,\beta_2)}
         & =T^{i_2j_2}_{\alpha_1\beta_1}
        (T^\dagger)^{i_1j_1}_{\alpha_2\beta_2}                                   \\
         & =(T^\dagger)^{i_1j_1}_{\alpha_2\beta_2}
        T^{i_2j_2}_{\alpha_1\beta_1}                                             \\
         & =[U_2]^{(i_1,i_2),(j_1,j_2)}_{(\alpha_2,\alpha_1),(\beta_2,\beta_1)}.
    \end{align}
    This is exactly the entrywise form of conjugation by $\mathbb S$.
\end{proof}

\begin{lemma}[Specified-tensor source-index values of the shifts]
    \label{thm:mpu_shift_specified_source_index_values}
    \lean{MPOTensor.sourceIndexValue_rightShiftTensor,
        MPOTensor.sourceIndexValue_leftShiftTensor,
        MPOTensor.sourceIndexValue_identityMPUTensor,
        MPOTensor.sourceIndexValue_shiftExampleU₁,
        MPOTensor.sourceIndexValue_shiftExampleU₂,
        MPOTensor.sourceIndexValue_shiftExampleU₃}
    \leanok
    \uses{def:mpu_specified_source_index_value,
        thm:mpu_shift_source_cuts_ranks,
        thm:mpu_tensor_product_source_ranks, threeMPU}
    Assume $d>0$. For the displayed right- and left-shift tensors $T$ and
    $T^\sharp$, the specified-tensor values are
    \begin{align}
        \operatorname{ind}_{\mathrm{src}}(T) & =\log_2 d.
    \end{align}
    Likewise,
    \begin{align}
        \operatorname{ind}_{\mathrm{src}}(T^\sharp) & =-\log_2 d.
    \end{align}
    The bond-one identity tensor and the three tensors $U_1,U_2,U_3$ satisfy
    \begin{align}
        \operatorname{ind}_{\mathrm{src}}(\Id)
        =\operatorname{ind}_{\mathrm{src}}(U_1)
        =\operatorname{ind}_{\mathrm{src}}(U_2)
        =\operatorname{ind}_{\mathrm{src}}(U_3) & =0.
    \end{align}
    Here $\operatorname{ind}_{\mathrm{src}}$ is the specified-tensor quantity
    from Definition~\ref{def:mpu_specified_source_index_value}, evaluated on the
    displayed unblocked tensors. It is not the public MPU index
    $\operatorname{ind}$ of Definition~\ref{def:index}, whose definition uses a
    simple blocking and requires independence of that choice.
    The public index values follow in
    Lemma~\ref{thm:mpu_shift_public_index_values}.
    % Source: paper_v2.tex lines 2037--2041.
\end{lemma}

\begin{proof}\leanok
    \uses{def:mpu_specified_source_index_value,
        thm:mpu_shift_source_cuts_ranks,
        thm:mpu_tensor_product_source_ranks, threeMPU}
    The right shift has source ranks $(d^2,1)$, while the left shift has
    $(1,d^2)$. Substitution into the source-index expression gives the first
    two formulas. The identity tensor has equal source ranks $(d,d)$. Source
    ranks multiply under the independent tensor product, so each of
    $U_1,U_2,U_3$ has equal right and left rank $d^2$. Their logarithmic
    differences therefore vanish.
\end{proof}

\begin{lemma}[Public indices of the shift examples]
    \label{thm:mpu_shift_public_index_values}
    \lean{MPOTensor.rightShiftTensor_exists_canonicalFormII,
        MPOTensor.rightShiftTensor_index, MPOTensor.leftShiftTensor_index,
        MPOTensor.shiftExampleU₁_index, MPOTensor.shiftExampleU₂_index,
        MPOTensor.shiftExampleU₃_index}
    \leanok
    \uses{def:mpu_canonical_form_ii, def:mpu_reduced_representative_index,
        threeMPU}
    Let $d>0$. The right shift is in canonical form II in its original bond
    coordinates, with ambient right fixed matrix $\rho=I/d$. The right and
    left shifts and the three two-spin families satisfy
    \begin{align}
        \operatorname{ind}(T)        & =\log_2 d,               &
        \operatorname{ind}(T^\sharp) & =-\log_2 d,                \\
        \operatorname{ind}(U_1)      & =\operatorname{ind}(U_2)
        =\operatorname{ind}(U_3)=0.
    \end{align}
    As asserted in \cite{Cirac2017MPU}, these are the blocking-independent indices of the families in
    Definition~\ref{threeMPU}.
    % Source: paper_v2.tex lines 2037--2041.
\end{lemma}

\begin{proof}\leanok
    \uses{thm:mpu_canonical_form_ii_stabilization_converse,
        thm:mpu_shift_blocked_source_index_values,
        prop:mpu_admissible_index_well_defined,
        thm:mpu_reduced_representative_index_invariance,
        thm:mpu_representative_index_blocking_adjoint,
        thm:mpu_representative_index_tensor_product}
    The normalized transfer map of the right shift sends $X$ to
    $\operatorname{tr}(X)I/d$. Its positive diagonal fixed matrix $I/d$
    gives a canonical-form-II presentation. The source-rank computation
    therefore yields the public index $\log_2 d$. Physical adjunction negates
    this index, giving the left-shift value. The identity has index zero,
    and additivity under the tensor product gives all three remaining values.
\end{proof}

\begin{lemma}[Source ranks and index values of the shifts at every blocking]
    \label{thm:mpu_shift_blocked_source_index_values}
    \lean{MPOTensor.evalWord_rightShiftTensor_ofFn,
        MPOTensor.blockTensor_rightShiftTensor_apply,
        MPOTensor.sourceCutM₁_blockTensor_rightShiftTensor,
        MPOTensor.sourceCutM₂_blockTensor_rightShiftTensor,
        MPOTensor.rightRank_blockTensor_rightShiftTensor,
        MPOTensor.leftRank_blockTensor_rightShiftTensor,
        MPOTensor.rightRank_mul_leftRank_blockTensor_rightShiftTensor,
        MPOTensor.sourceIndexValue_blockTensor_rightShiftTensor,
        MPOTensor.blockTensor_leftShiftTensor,
        MPOTensor.rightRank_blockTensor_leftShiftTensor,
        MPOTensor.leftRank_blockTensor_leftShiftTensor,
        MPOTensor.sourceIndexValue_blockTensor_leftShiftTensor,
        MPOTensor.sourceRanks_blockTensor_shiftExampleU₂,
        MPOTensor.sourceRanks_blockTensor_shiftExampleU₃,
        MPOTensor.sourceIndexValue_blockTensor_shiftExampleU₂,
        MPOTensor.sourceIndexValue_blockTensor_shiftExampleU₃,
        MPOTensor.rightShiftTensor_isMPUSimple,
        MPOTensor.leftShiftTensor_isMPUSimple,
        MPOTensor.blockTensor_rightShiftTensor_isMPUSimple,
        MPOTensor.blockTensor_leftShiftTensor_isMPUSimple,
        MPOTensor.identityMPUTensor_isMPUSimple,
        MPOTensor.shiftExampleU₁_isMPUSimple,
        MPOTensor.shiftExampleU₂_isMPUSimple,
        MPOTensor.shiftExampleU₃_isMPUSimple,
        MPOTensor.blockTensor_shiftExampleU₂_isMPUSimple,
        MPOTensor.blockTensor_shiftExampleU₃_isMPUSimple,
        MPOTensor.rightRank_blockTensor_rightShiftTensor_ne_leftRank,
        MPOTensor.sourceIndexValue_blockTensor_rightShiftTensor_ne_zero}
    \leanok
    \uses{def:mpu_simple, thm:mpu_bounded_simple_blocking,
        def:mpu_specified_source_index_value,
        thm:mpu_specified_source_index_identities,
        thm:mpu_physical_adjoint_source_ranks,
        thm:mpu_tensor_product_source_ranks, thm:mpu_source_ranks_physical_reindexing,
        threeMPU, thm:mpu_block_tensor_product, thm:mpu_simple_tensor_product}
    Assume $d>0$ and let $k\ge0$. The right-shift tensor $T$ and the
    left-shift tensor $T^\sharp$ are simple, and so are their blockings
    $T_{k+1}$ and $T^\sharp_{k+1}$ over $k+1$ sites. The blocked letters are
    \begin{align}
        T_{k+1}^{\sigma\tau}
        =\Bigl(\prod_{m<k}\delta_{\tau_m,\sigma_{m+1}}\Bigr)\,
        |\sigma_0)(\tau_k|,
    \end{align}
    and the source ranks are
    \begin{align}
        (r,\ell)(T_{k+1})        & =(d^{k+2},d^{k}), &
        (r,\ell)(T^\sharp_{k+1}) & =(d^{k},d^{k+2}).
    \end{align}
    Both source ranks of the blockings of $U_2$ and $U_3$ equal $d^{2k+2}$.
    Consequently
    \begin{align}
        \operatorname{ind}_{\mathrm{src}}(T_{k+1})        & =\log_2 d,  &
        \operatorname{ind}_{\mathrm{src}}(T^\sharp_{k+1}) & =-\log_2 d,
    \end{align}
    and the blockings of $U_2$ and $U_3$ have value zero. If $d>1$, then
    $r(T_{k+1})\ne\ell(T_{k+1})$ and
    $\operatorname{ind}_{\mathrm{src}}(T_{k+1})\ne0$.

    The tensors $U_1$, $U_2$, and $U_3$ are simple, and so are the blockings
    of $U_2$ and $U_3$ over $k+1$ sites, because independent tensor products
    of simple tensors are simple. Thus all four values are taken at simple
    blockings, which is the simplicity requirement of Definition~\ref{def:index}.
    The canonical-form construction and representative invariance in
    Lemma~\ref{thm:mpu_shift_public_index_values} identify these values with
    the public indices and carry them over to any positive-dimensional MPU
    generating the same periodic operators. Unequal ranks do not by
    themselves exclude a circuit: after blocking, every matrix product unitary,
    the shift included, is a depth-two circuit of nearest-neighbour unitaries
    in its standard form. The statement that the shift is not a finite-depth
    circuit of local gates needs parts (ii)--(iv) of Theorem~\ref{IndexTh}
    together with index zero for one layer of nearest-neighbour gates, and is
    not asserted here.
    % The former specified-tensor restriction is resolved in
    % docs/paper-gaps/mpu_shift_specified_tensor_index_scope.tex.
    % Source: paper_v2.tex lines 605--622, 681--704, 1980--2001, 2037--2041, 2308;
    % review TN-Review-main.tex lines 2607--2611.
\end{lemma}

\begin{proof}\leanok
    \uses{def:mpu_simple, thm:mpu_bounded_simple_blocking,
        thm:mpu_physical_adjoint_source_ranks,
        thm:mpu_tensor_product_source_ranks,
        thm:mpu_source_ranks_physical_reindexing,
        thm:mpu_specified_source_index_identities, thm:mpu_block_tensor_product, thm:mpu_double_layer_blocking,
        thm:mpu_simple_tensor_product}
    With $\Phi=\sum_j|jj)$, the double-layer letters of $T$ are
    $|\Phi)(ik|$, so the boundary vectors $|00)$ and $\Phi$ satisfy both
    simplicity identities; for $T^\sharp$ the letters are $|ik)(\Phi|$ and the
    boundary vectors are exchanged. The identity tensor has bond dimension
    one and double-layer letters $\delta_{ij}$, so $a=b=1$ are boundary
    vectors. Simplicity survives independent tensor products and blocking,
    which gives simplicity of $U_1$, $U_2$, $U_3$ and of the blockings of
    $U_2$ and $U_3$. A product
    of matrix units $|\sigma_m)(\tau_m|$ is $|\sigma_0)(\tau_k|$ when adjacent
    indices match and zero otherwise, which gives the blocked letters. The
    first source cut of $T_{k+1}$ sends the row $(\sigma,\beta)$ to the single
    column $(\sigma_0,\ \sigma_1\cdots\sigma_k\beta)$; this is a bijection, so
    the cut is a permutation matrix of rank $d^{k+2}$. The second source cut
    factors through the $d^k$ middle words $\sigma_1\cdots\sigma_k$, and its
    rows $(0,0y)$ and columns $(y0,0)$ form an identity minor of size $d^k$.
    Physical adjunction exchanges the two ranks, and ranks multiply under
    tensor products and are unchanged by physical reindexing. Substituting
    into the source-index expression gives
    $\frac12\bigl((k+2)-k\bigr)\log_2 d=\log_2 d$ and the remaining values.
    For $d>1$ the powers $d^{k+2}$ and $d^k$ differ.
\end{proof}

\begin{definition}[Explicit source factors for the three shift MPUs]
    \label{def:threeMPU_supplied_source_factors}
    \lean{MPOTensor.identitySourceFactors,
        MPOTensor.rightShiftSourceFactors,
        MPOTensor.leftShiftSourceFactors,
        MPOTensor.SourceFactors.independentTensorProductOfIdentityWeight,
        MPOTensor.identityTensorIdentityMatrix,
        MPOTensor.identityTensorIdentityMatrix_apply,
        MPOTensor.swapTensorSwapMatrix_mul_identitySwapIdentityMatrix_apply,
        MPOTensor.identitySwapIdentityMatrix_mul_swapTensorSwapMatrix_apply,
        MPOTensor.shiftExampleU₁LeftRankEquiv,
        MPOTensor.shiftExampleU₁RightRankEquiv,
        MPOTensor.shiftExampleU₂LeftRankEquiv,
        MPOTensor.shiftExampleU₂RightRankEquiv,
        MPOTensor.shiftExampleU₃RightRankEquiv_apply,
        MPOTensor.shiftExampleU₁SourceFactors,
        MPOTensor.shiftExampleU₂SourceFactors,
        MPOTensor.shiftExampleU₃SourceFactors}
    \leanok
    \uses{threeMPU, def:mpu_weighted_source_factors, thm:mpu_shift_source_cuts_ranks}
    Assume $d>0$. Equip $U_1$ with the tensor product of the two bond-one
    identity source factorizations.  Equip $U_2$ with the tensor product of the
    left-shift and right-shift source factorizations, in that order, and equip
    $U_3$ with the tensor product in the reverse order.  In all three cases the
    virtual weight is the identity.  These are explicit source
    factorizations; no identification with an arbitrary compact singular-value
    decomposition is asserted.
\end{definition}

\begin{lemma}[Unitarity of the four-spin shift matrices]
    \label{thm:threeMPU_four_spin_matrix_unitarity}
    \lean{MPOTensor.identitySwapIdentityMatrix_isUnitaryBetween,
        MPOTensor.swapTensorSwapMatrix_isUnitaryBetween,
        MPOTensor.swapTensorSwapMatrix_mul_identitySwapIdentityMatrix_isUnitaryBetween,
        MPOTensor.identitySwapIdentityMatrix_mul_swapTensorSwapMatrix_isUnitaryBetween}
    \leanok
    Let $\mathbb S$ exchange two $d$-dimensional spins. The two swap factors
    appearing in the blocked forms of $U_2$ and $U_3$, and both ordered
    products, are unitary:
    \begin{align}
        \Id\otimes\mathbb S\otimes\Id,\quad
        \mathbb S\otimes\mathbb S,\quad
        (\mathbb S\otimes\mathbb S)(\Id\otimes\mathbb S\otimes\Id),\quad
        (\Id\otimes\mathbb S\otimes\Id)(\mathbb S\otimes\mathbb S).
    \end{align}
    % Source lines: paper_v2.tex 2018--2034.
\end{lemma}

\begin{proof}\leanok
    \uses{thm:mpu_swap_matrix_unitary}
    The swap is unitary by Lemma~\ref{thm:mpu_swap_matrix_unitary}. Kronecker
    products preserve isometries, and a square isometry is unitary. Thus both
    factors are unitary. Closure of the unitary matrices under multiplication
    gives the two ordered products.
\end{proof}

\begin{lemma}[Explicit paper gates for $U_1$]
    \label{thm:threeMPU_supplied_source_gates_one}
    \lean{MPOTensor.shiftExampleU₁SourceURowEquiv,
        MPOTensor.shiftExampleU₁SourceVColumnEquiv,
        MPOTensor.shiftExampleU₁_sourceU_fourSpin_apply,
        MPOTensor.shiftExampleU₁_sourceV_fourSpin_apply}
    \leanok
    \uses{def:threeMPU_supplied_source_factors, def:mpu_source_uv}
    In the four-spin coordinates of the explicit source factorization,
    \begin{align}
        u_1 & =\Id\otimes\Id,
            & v_1             & =\Id\otimes\Id.
    \end{align}
\end{lemma}

\begin{proof}\leanok
    \uses{def:threeMPU_supplied_source_factors, def:mpu_source_uv}
    The tensor-product formulas for $u=Y_2\mathbin{-}Y_1$ and
    $v=X_1\mathbin{-}X_2$ reduce each entry to the product of two primitive
    identity entries.
\end{proof}

\begin{lemma}[Paper gates in the two-site factorizations of $U_2$]
    \label{thm:threeMPU_U2_supplied_source_gates_two}
    \lean{MPOTensor.shiftExampleU₂SourceURowEquiv,
        MPOTensor.shiftExampleU₂SourceVColumnEquiv,
        MPOTensor.shiftExampleU₂_sourceU_fourSpin_apply,
        MPOTensor.shiftExampleU₂_sourceV_fourSpin_apply,
        MPOTensor.shiftExampleU₂_blockTwo_apply_eq_sum_X₂_mul_sourceU_mul_X₁,
        MPOTensor.shiftExampleU₂_blockTwo_apply_eq_sum_sourceV_mul_Y₁_mul_Y₂}
    \leanok
    \uses{def:threeMPU_supplied_source_factors, def:mpu_source_uv,
        thm:mpu_source_v_open_two_site, thm:mpu_source_u_open_two_site}
    Assume $d>0$. In the four-spin coordinates of the paper, the explicit
    source gates satisfy
    \begin{align}
        u_2^{(2)}    & =d\,(\Id\otimes\mathbb S\otimes\Id), \\
        d\,v_2^{(2)} & =(\mathbb S\otimes\mathbb S)
        (\Id\otimes\mathbb S\otimes\Id).
    \end{align}
    The blocked tensor factors through $X_2,u_2^{(2)},X_1$ and, equivalently,
    through $v_2^{(2)},Y_1,Y_2$. The scalar factors record the normalization of
    the explicit source factorizations.
\end{lemma}

\begin{proof}\leanok
    \uses{def:threeMPU_supplied_source_factors, def:mpu_source_uv,
        thm:mpu_source_v_open_two_site, thm:mpu_source_u_open_two_site}
    The tensor-product formulas for the paper gates reduce their entries to
    products of the left- and right-shift source gates. The Kronecker deltas
    give the two displayed swap matrices, with the factors $d$ and $d^{-1}$
    fixed by the identity-weight normalization; substituting these evaluations
    into the two exact open source-gate factorizations yields the displayed gates.
\end{proof}

\begin{lemma}[Paper gates in the two-site factorizations of $U_3$]
    \label{thm:threeMPU_U3_supplied_source_gates_two}
    \lean{MPOTensor.shiftExampleU₃SourceURowEquiv,
        MPOTensor.shiftExampleU₃SourceVColumnEquiv,
        MPOTensor.shiftExampleU₃_sourceU_fourSpin_apply,
        MPOTensor.shiftExampleU₃_sourceV_fourSpin_apply,
        MPOTensor.shiftExampleU₃_blockTwo_apply_eq_sum_X₂_mul_sourceU_mul_X₁,
        MPOTensor.shiftExampleU₃_blockTwo_apply_eq_sum_sourceV_mul_Y₁_mul_Y₂}
    \leanok
    \uses{def:threeMPU_supplied_source_factors, def:mpu_source_uv,
        thm:mpu_source_v_open_two_site, thm:mpu_source_u_open_two_site}
    Assume $d>0$. In the four-spin coordinates of the paper, the explicit
    source gates satisfy
    \begin{align}
        u_3^{(2)}    & =d\,(\Id\otimes\mathbb S\otimes\Id)
        (\mathbb S\otimes\mathbb S),                       \\
        d\,v_3^{(2)} & =\Id\otimes\mathbb S\otimes\Id.
    \end{align}
    The blocked tensor factors through $X_2,u_3^{(2)},X_1$ and, equivalently,
    through $v_3^{(2)},Y_1,Y_2$. The scalar factors record the normalization of
    the explicit source factorizations.
\end{lemma}

\begin{proof}\leanok
    \uses{def:threeMPU_supplied_source_factors, def:mpu_source_uv,
        thm:mpu_source_v_open_two_site, thm:mpu_source_u_open_two_site}
    Reversing the two primitive shift factors in the tensor-product formulas for
    $u$ and $v$ makes their Kronecker deltas give the displayed ordered swap products,
    with the factors $d$ and $d^{-1}$ fixed by the identity-weight normalization, and
    substituting these evaluations into the two exact open source-gate
    factorizations completes the calculation.
\end{proof}

\begin{lemma}[Trace-normalized source gates for the counterpropagating shifts]
    \label{thm:threeMPU_normalized_source_gates}
    \lean{MPOTensor.shiftExampleU₂Paper_sourceU_fourSpin_apply,
        MPOTensor.shiftExampleU₂Paper_sourceV_fourSpin_apply,
        MPOTensor.shiftExampleU₃Paper_sourceU_fourSpin_apply,
        MPOTensor.shiftExampleU₃Paper_sourceV_fourSpin_apply,
        MPOTensor.shiftExampleU₂PaperSourceFactors,
        MPOTensor.shiftExampleU₃PaperSourceFactors}
    \leanok
    \uses{threeMPU, def:mpu_weighted_source_factors,
        def:mpu_source_uv, def:mpu_shift_trace_weight}
    For $d>0$, the two counterpropagating shift tensors have source factors
    with trace-one weight $\rho=d^{-2}\Id_{d^2}$. In the four-spin coordinates
    used above, their source gates are
    \begin{align}
        u_2^{(2)} & = \Id\otimes\mathbb S\otimes \Id,
                  & v_2^{(2)}                                                      & = (\mathbb S\otimes\mathbb S)(\Id\otimes\mathbb S\otimes \Id), \\
        u_3^{(2)} & = (\Id\otimes\mathbb S\otimes \Id)(\mathbb S\otimes\mathbb S),
                  & v_3^{(2)}                                                      & = \Id\otimes\mathbb S\otimes \Id.
    \end{align}
\end{lemma}
\begin{proof}\leanok
    \uses{def:threeMPU_supplied_source_factors, thm:threeMPU_U2_supplied_source_gates_two,
        thm:threeMPU_U3_supplied_source_gates_two}
    Starting with the identity-weight source factors, replace
    $X_1,Y_1,Z_1$ by $dX_1,d^{-1}Y_1,dZ_1$, respectively, and leave the
    second cut unchanged. The cut factorizations and right inverses are
    preserved, while the first weighted isometry now uses $d^{-2}\Id_{d^2}$.
    Writing $\hat u_k,\hat v_k$ for the identity-weight gates of
    Lemma~\ref{thm:threeMPU_U2_supplied_source_gates_two} and
    Lemma~\ref{thm:threeMPU_U3_supplied_source_gates_two}, the source gates
    rescale as
    \begin{align}
        u_k^{(2)} & = d^{-1}\hat u_k, & v_k^{(2)} & = d\,\hat v_k, & k & = 2,3,
    \end{align}
    which cancels the scalar factors in those formulas.
\end{proof}

\begin{lemma}[Two-site decompositions of the three shift MPUs]
    \label{thm:threeMPU_supplied_decompositions}
    \lean{MPOTensor.shiftExampleU₁StandardForm_normalized,
        MPOTensor.shiftExampleU₃StandardForm_normalized,
        MPOTensor.shiftExampleU₂BlockedStandardForm_normalized,
        MPOTensor.shiftExampleU₃BlockedStandardForm_normalized,
        MPOTensor.shiftExampleU₁StandardForm,
        MPOTensor.shiftExampleU₃StandardForm,
        MPOTensor.shiftExampleU₂BlockedStandardForm,
        MPOTensor.shiftExampleU₃BlockedStandardForm}
    \leanok
    \uses{threeMPU, def:mpu_supplied_two_site_standard_form,
        def:mpu_shift_trace_weight}
    Assume $d>0$. The three shift families have two-site decompositions
    with given gates whose unitary gates satisfy
    \begin{align}
        u_1                          & =\Id\otimes\Id,                  & v_1 & =\Id\otimes\Id, \\
        u_3                          & =\mathbb S,                      & v_3 & =\mathbb S,     \\
        u_2^{(2)}                    & =\Id\otimes\mathbb S\otimes\Id,  &
        v_2^{(2)}                    & =(\mathbb S\otimes\mathbb S)
        (\Id\otimes\mathbb S\otimes\Id),                                                        \\
        u_3^{(2)}                    & =(\Id\otimes\mathbb S\otimes\Id)
        (\mathbb S\otimes\mathbb S), &
        v_3^{(2)}                    & =\Id\otimes\mathbb S\otimes\Id.
    \end{align}
    The two-site blockings in the last two lines use the physical order
    $((1,2),(3,4))$. These are unitary gates with normalized half factors
    satisfying the open standard-form equation. For the two blocked families,
    the half factors and gates are those of the source factors of
    Lemma~\ref{thm:threeMPU_normalized_source_gates}. For $U_1$ and the
    unblocked $U_3$, the one-site tensor is read as a pair of spins and
    decomposed directly; as recorded in \cite{gap:mpu_shift_supplied_standard_forms}, these decompositions are not identified with the
    standard form of
    Definition~\ref{SF}.
    % Source lines: paper_v2.tex 2002--2042.
\end{lemma}

\begin{proof}\leanok
    \uses{def:mpu_source_weight, thm:threeMPU_normalized_source_gates,
        thm:mpu_source_u_open_two_site, thm:threeMPU_four_spin_matrix_unitarity,
        thm:mpu_swap_matrix_unitary}
    For the identity family, take
    $X_1(i,\gamma;r)=\delta_{i,r}$ and
    $X_2(\alpha,i;\ell)=\delta_{i,\ell}$. For the unblocked third family, take
    \begin{align}
        X_1(i,(a,b);r)    & =\sqrt d\,\delta_{a,r}\delta_{b,i},   &
        X_2((a,b),i;\ell) & =d^{-1/2}\delta_{a,i}\delta_{b,\ell}.
    \end{align}
    Its gates are the swap, which is unitary by
    Lemma~\ref{thm:mpu_swap_matrix_unitary}.
    The reciprocal scalars cancel in both contractions. Summing the Gram
    entries proves $X_1^*(\Id\otimes\rho)X_1=\Id$ and $X_2^*X_2=\Id$, with
    $\rho=d^{-2}\Id_{d^2}$ for the third family and $\rho=1$ for the identity.
    For the two blocked families, the trace-normalized source factors
    of Lemma~\ref{thm:threeMPU_normalized_source_gates} apply after relabeling the
    intermediate coordinates by pairs of spin indices. With
    $e_L:\{0,\ldots,d^2-1\}\to L$ and
    $e_R:\{0,\ldots,d^2-1\}\to R$ denoting these bijections, obtained by
    $\ell\leftrightarrow(a,b)$ and $r\leftrightarrow(c,e)$ and the
    source-factor rank identifications. The open contraction becomes
    \begin{align}
        \sum_{p\in L}\sum_{q\in R}
        X_2(\alpha,i_1;p)u(p,q;j_1,j_2)X_1(i_2,\gamma;q)
         & =\sum_{\ell=0}^{d^2-1}\sum_{r=0}^{d^2-1}
        X_2(\alpha,i_1;e_L(\ell))u(e_L(\ell),e_R(r);j_1,j_2)
        X_1(i_2,\gamma;e_R(r)).
        \notag
    \end{align}
    Thus the source-factor contraction theorem gives the open equation
    in the displayed coordinates. The four-spin permutation matrices are
    unitary. The relabeled half factors are
    $\widehat X_1(i,\gamma;r)=X_1(i,\gamma;e_R(r))$ and
    $\widehat X_2(\alpha,i;\ell)=X_2(\alpha,i;e_L(\ell))$. Since $e_R$
    and $e_L$ are injective, the source normalization identities give
    \begin{align}
        \bigl(\widehat X_1^*(\Id\otimes\rho)\widehat X_1\bigr)_{r,t}
         & =\bigl(X_1^*(\Id\otimes\rho)X_1\bigr)_{e_R(r),e_R(t)}
        =\delta_{e_R(r),e_R(t)}=\delta_{r,t},                    \\
        \bigl(\widehat X_2^*\widehat X_2\bigr)_{\ell,m}
         & =\bigl(X_2^*X_2\bigr)_{e_L(\ell),e_L(m)}
        =\delta_{e_L(\ell),e_L(m)}=\delta_{\ell,m}.
    \end{align}
\end{proof}

\begin{definition}[Transformation with swaps]
    \label{threeMPU2}
    \lean{MPOTensor.shiftPhysicalSwap}
    \lean{MPOTensor.shiftPhysicalSwap_apply}
    \lean{MPOTensor.shiftExampleTildeU₁}
    \lean{MPOTensor.shiftExampleTildeU₂}
    \lean{MPOTensor.shiftExampleTildeU₃}
    \leanok
    \uses{threeMPU, def:mpu_physical_pair_swap}
    Let $\mathbb S$ exchange the two $d$-dimensional spins at one site, with
    matrix entries
    \begin{align}
        \mathbb S_{(p,q),(r,s)}=\delta_{p,s}\delta_{q,r}.
    \end{align}
    With the row index as the output physical index and the column index as the
    input physical index, define
    \begin{align}
        \widetilde{\mathcal U}_a^{\,IJ}
        =\sum_K\mathbb S_{IK}\mathcal U_a^{\,KJ},
        \qquad a\in\{1,2,3\}.
    \end{align}
    Thus the swap acts on the output leg, and the corresponding closed operator
    is left-multiplied by $\mathbb S^{\otimes N}$.
    % Source lines: paper_v2.tex 2045--2062.
\end{definition}

\begin{lemma}[Finite-chain formulas for the swap-transformed shifts]
    \label{thm:threeMPU2_chain_formulas}
    \lean{MPOTensor.mpo_ketLeftMul}
    \lean{MPOTensor.IsMPU.ketLeftMul}
    \lean{MPOTensor.shiftPhysicalSwap_mem_unitaryGroup}
    \lean{MPOTensor.mpo_shiftExampleTildeU₁}
    \lean{MPOTensor.mpo_shiftExampleTildeU₂}
    \lean{MPOTensor.mpo_shiftExampleTildeU₃}
    \lean{MPOTensor.shiftExampleTildeU₁_isMPU}
    \lean{MPOTensor.shiftExampleTildeU₂_isMPU}
    \lean{MPOTensor.shiftExampleTildeU₃_isMPU}
    \leanok
    \uses{threeMPU2, thm:threeMPU_chain_formulas,
        def:mpdo_sitewise_physical_matrix, def:mpu}
    On every nonempty chain,
    \begin{align}
        \widetilde U_1^{(N)} & =\mathbb S^{\otimes N}, \\
        \widetilde U_2^{(N)} & =\mathbb S^{\otimes N}
        (T^{(N)\dagger}\otimes T^{(N)}),               \\
        \widetilde U_3^{(N)} & =\mathbb S^{\otimes N}
        (T^{(N)}\otimes T^{(N)\dagger}).
    \end{align}
    Each transformed tensor is a matrix product unitary.
    % Source lines: paper_v2.tex 2052--2062.
\end{lemma}

\begin{proof}\leanok
    \uses{threeMPU2, thm:threeMPU_chain_formulas,
        def:mpdo_sitewise_physical_matrix, def:mpu}
    Acting by a one-site matrix $P$ on every output physical index gives
    \begin{align}
        \operatorname{MPO}_N(P\mathcal M)
        =P^{\otimes N}\operatorname{MPO}_N(\mathcal M).
    \end{align}
    Substituting $P=\mathbb S$ and the three shift formulas gives the stated
    identities. The swap is a permutation matrix and hence unitary. If
    $\mathcal M$ is an MPU, then
    \begin{align}
        \bigl(P^{\otimes N}\operatorname{MPO}_N(\mathcal M)\bigr)^\dagger
        \bigl(P^{\otimes N}\operatorname{MPO}_N(\mathcal M)\bigr)
         & =\operatorname{MPO}_N(\mathcal M)^\dagger
        (P^{\otimes N})^\dagger P^{\otimes N}
        \operatorname{MPO}_N(\mathcal M)             \\
         & =\Id.
    \end{align}
    Thus the sitewise left action preserves the MPU property.
\end{proof}

\begin{lemma}[Conjugacy of the third swap-transformed shift]
    \label{thm:tildeU3}
    \lean{MPOTensor.mpo_shiftExampleTildeU₃_eq_swap_mul_tildeU₂_mul_swap}
    \leanok
    \uses{threeMPU2, thm:threeMPU2_chain_formulas}
    On every nonempty chain,
    \begin{align}
        \widetilde U_3^{(N)}
        =\mathbb S^{\otimes N}\widetilde U_2^{(N)}\mathbb S^{\otimes N}.
    \end{align}
    In \cite[lines~2059--2062]{Cirac2017MPU}, this equation is the swap-conjugacy
    identity.
    % Source label: tildeU3.
\end{lemma}

\begin{proof}\leanok
    \uses{thm:threeMPU2_chain_formulas}
    Splitting each physical configuration into its two spin configurations represents the
    three operators as permutation matrices.  In these coordinates,
    $\mathbb S^{\otimes N}$ exchanges the two configurations,
    $U_2^{(N)}$ acts by $(T^{(N)\dagger},T^{(N)})$, and $U_3^{(N)}$ acts by
    $(T^{(N)},T^{(N)\dagger})$.  Exchanging the two factors on both sides of
    the second action gives the third, which proves the identity.
\end{proof}

\begin{lemma}[Two-site decompositions of the swap-transformed shifts]
    \label{thm:threeMPU2_supplied_decompositions}
    \lean{MPOTensor.shiftExampleTildeU₁StandardForm_normalized,
        MPOTensor.shiftExampleTildeU₂StandardForm_normalized,
        MPOTensor.shiftExampleTildeU₁StandardForm,
        MPOTensor.shiftExampleTildeU₂StandardForm}
    \leanok
    \uses{threeMPU2, def:mpu_supplied_two_site_standard_form,
        def:mpu_shift_trace_weight}
    For $d>0$, the first two transformed families have two-site
    decompositions with given gates and normalized half factors, with gates
    \begin{align}
        \widetilde u_1 & =\mathbb S,     &
        \widetilde v_1 & =\Id\otimes\Id,   \\
        \widetilde u_2 & =\Id\otimes\Id, &
        \widetilde v_2 & =\mathbb S.
    \end{align}
    The one-site tensors are read as pairs of spins and decomposed directly;
    as recorded in \cite{gap:mpu_shift_supplied_standard_forms}, these decompositions are not identified with the standard form of
    Definition~\ref{SF}.
    % Source lines: paper_v2.tex 2087--2099.
\end{lemma}

\begin{proof}\leanok
    \uses{threeMPU2, thm:threeMPU_supplied_decompositions,
        def:mpu_source_weight, def:mpu_supplied_two_site_standard_form,
        thm:mpu_swap_matrix_unitary}
    The physical swap exchanges the two output indices of each tensor.
    For the first family, retain the identity family's half factors and
    replace its unshifted gate by the swap. For the second family, choose
    $X_1(i,(a,b);r)=\sqrt d\,\delta_{a,i}\delta_{b,r}$ and
    $X_2((a,b),i;\ell)=d^{-1/2}\delta_{a,\ell}\delta_{b,i}$.
    Contracting over the bond index $\beta=(a,b)$ gives the shifted swap
    gate,
    \begin{align}
        \sum_{a,b}X_1(i_1,(a,b);s)\,X_2((a,b),i_2;t)
        =\delta_{i_1,t}\delta_{i_2,s}
        =\mathbb S_{(i_1,i_2),(s,t)}.
    \end{align}
    Inserting the unshifted identity gate $\Id\otimes\Id$ between the half
    factors gives, for $\alpha=(a,b)$ and $\gamma=(a',b')$,
    \begin{align}
        \sum_{\ell,r}X_2(\alpha,i_1;\ell)\,
        (\Id\otimes\Id)_{(\ell,r),(j_1,j_2)}\,X_1(i_2,\gamma;r)
        =X_2(\alpha,i_1;j_1)\,X_1(i_2,\gamma;j_2)
        =\delta_{a,j_1}\delta_{b,i_1}\delta_{a',i_2}\delta_{b',j_2},
    \end{align}
    which is the entry of $\widetilde{\mathcal U}_2$ with output
    $(i_1,i_2)$, input $(j_1,j_2)$, and bond indices $\alpha,\gamma$.
    The swap is unitary by Lemma~\ref{thm:mpu_swap_matrix_unitary}, so
    both gates are unitary. With $\rho=d^{-2}\Id_{d^2}$, the Gram sums are
    \begin{align}
        \bigl(X_1^*(\Id\otimes\rho)X_1\bigr)_{r,r'}
         & =d^{-2}\sum_{i,a,b}d\,\delta_{a,i}\delta_{b,r}\delta_{b,r'}
        =\delta_{r,r'},                                                   \\
        \bigl(X_2^*X_2\bigr)_{\ell,\ell'}
         & =d^{-1}\sum_{a,b,i}\delta_{a,\ell}\delta_{a,\ell'}\delta_{b,i}
        =\delta_{\ell,\ell'},
    \end{align}
    since $d$ pairs $(i,a)$ satisfy $a=i$ and $d$ pairs $(b,i)$ satisfy $b=i$.
\end{proof}

\begin{definition}[Adjunction, transposition, and conjugation symmetries]
    \label{def:mpu_swap_comparison_symmetries}
    \lean{MPOTensor.daggerSymmetry,
        MPOTensor.transposeSymmetry,
        MPOTensor.conjugationSymmetry,
        MPOTensor.shiftSwapDaggerSymmetry,
        MPOTensor.shiftSwapTransposeSymmetry}
    \leanok
    \uses{def:mpu_finite_chain_operator_symmetry, threeMPU2}
    On every positive-length chain, the adjunction, transposition, and
    conjugation actions are $U\mapsto U^\dagger$, $U\mapsto U^{\mathsf T}$,
    and $U\mapsto\overline U$, respectively. For a site consisting of two
    $d$-dimensional spins, let $S_N=\mathbb S^{\otimes N}$ be their sitewise
    exchange. The swap-combined actions are
    $U\mapsto S_NU^\dagger S_N$ and
    $U\mapsto S_NU^{\mathsf T}S_N$.
    These are the symmetries of
    \cite[Sections~IV.A--IV.C]{Cirac2017MPU}.
\end{definition}

\begin{theorem}[Pointwise comparison of symmetry invariance]
    \label{thm:mpu_pointwise_symmetry_invariance_comparison}
    \lean{MPOTensor.isInvariantUnderSymmetry_iff_of_pointwise}
    \leanok
    \uses{def:mpu_finite_chain_operator_symmetry}
    Let $\mathcal S,\mathcal T$ have the same applicable chain lengths,
    and let $\mathcal U,\mathcal V$ have the same physical and virtual
    dimensions. Suppose that, for every $N\geq0$,
    $\mathcal S[U^{(N)}]=U^{(N)}$ if and only if
    $\mathcal T[V^{(N)}]=V^{(N)}$. Then $\mathcal U$ is
    $\mathcal S$-invariant if and only if $\mathcal V$ is
    $\mathcal T$-invariant.
\end{theorem}

\begin{proof}\leanok
    \uses{def:mpu_finite_chain_operator_symmetry}
    Restrict the pointwise equivalences to the common set of applicable
    chain lengths and quantify over that set.
\end{proof}

\begin{theorem}[An involution converts a combined fixed equation]
    \label{thm:mpu_involution_combined_fixed_equation}
    \lean{Matrix.involution_conj_eq_iff_leftMul}
    \leanok
    Let $Q,X$ be complex matrices indexed by the same finite set, and let
    $F$ be a map on that matrix space. If $Q^2=I$ and
    $F(QX)=F(X)Q$, then
    \begin{align}
        QF(X)Q=X
         & \quad\Longleftrightarrow\quad F(QX)=QX.
        \label{eq:mpu_involution_combined_fixed}
    \end{align}
\end{theorem}

\begin{proof}\leanok
    Replace $F(QX)$ by $F(X)Q$ and multiply by $Q$ on the left.
    Involutivity gives both implications in
    (\ref{eq:mpu_involution_combined_fixed}).
\end{proof}

\begin{theorem}[Symmetry invariance under the one-sided swap]
    \label{thm:mpu_one_sided_swap_invariance}
    \lean{MPOTensor.shiftPhysicalSwap_mul_self,
        MPOTensor.isInvariantUnderSymmetry_shiftSwapDagger_iff,
        MPOTensor.isInvariantUnderSymmetry_shiftSwapTranspose_iff,
        MPOTensor.isInvariantUnderSymmetry_conjugation_ketLeftMul_swap_iff}
    \leanok
    \uses{def:mpu_swap_comparison_symmetries, threeMPU2}
    Let $\mathcal U$ be a local tensor with two $d$-dimensional physical
    spins, and let $\widetilde{\mathcal U}=\mathbb S\mathcal U$ denote the
    action on its output physical leg. The swap satisfies $\mathbb S^2=I$.
    For every $N>0$,
    $\widetilde U^{(N)}=S_NU^{(N)}$. Then $\mathcal U$ is invariant under
    swap-combined adjunction if and only if $\widetilde{\mathcal U}$ is
    invariant under adjunction, $\mathcal U$ is invariant under
    swap-combined transposition if and only if $\widetilde{\mathcal U}$ is
    invariant under transposition, and $\mathcal U$ is invariant under
    conjugation if and only if $\widetilde{\mathcal U}$ is.
    The one-sided transformation is the one prescribed by the transformed
    examples in \cite[Section~VI.B.1]{Cirac2017MPU}; the printed proof's
    two-sided expression has the local correction recorded in
    \cite{gap:mpu_swap_symmetry_one_sided}.
\end{theorem}

\begin{proof}\leanok
    \uses{thm:threeMPU2_chain_formulas,
        thm:mpu_pointwise_symmetry_invariance_comparison,
        thm:mpu_involution_combined_fixed_equation}
    The swap satisfies
    $S_N^2=I$, $S_N^\dagger=S_N^{\mathsf T}=S_N$, and
    $\overline{S_N}=S_N$. We show that, for every $N>0$,
    \begin{align}
        S_N(U^{(N)})^\dagger S_N=U^{(N)}
         & \quad\Longleftrightarrow\quad
        (\widetilde U^{(N)})^\dagger=\widetilde U^{(N)},
        \label{eq:mpu_swap_transport_dagger}    \\
        S_N(U^{(N)})^{\mathsf T}S_N=U^{(N)}
         & \quad\Longleftrightarrow\quad
        (\widetilde U^{(N)})^{\mathsf T}=\widetilde U^{(N)},
        \label{eq:mpu_swap_transport_transpose} \\
        \overline{U^{(N)}}=U^{(N)}
         & \quad\Longleftrightarrow\quad
        \overline{\widetilde U^{(N)}}=\widetilde U^{(N)}.
        \label{eq:mpu_swap_transport_conjugation}
    \end{align}
    Substituting $\widetilde U^{(N)}=S_NU^{(N)}$ and multiplying by
    $S_N$ proves (\ref{eq:mpu_swap_transport_dagger}) and
    (\ref{eq:mpu_swap_transport_transpose}). Reality of $S_N$ proves
    (\ref{eq:mpu_swap_transport_conjugation}). Since all three symmetries
    apply at the same chain lengths, the pointwise equivalences give the
    three equivalences of invariance.
\end{proof}

\begin{theorem}[Continuous involutive physical left actions]
    \label{thm:mpu_continuous_involutive_left_action}
    \lean{MPOTensor.continuous_ketLeftMul,
        MPOTensor.ketLeftMul_involutive}
    \leanok
    \uses{def:mpdo_vertical_one_site_actions}
    Fix physical and virtual dimensions $d,D$, and a matrix $Q\in\MN{d}$.
    The map $\mathcal U\mapsto Q\mathcal U$, acting on the output physical
    leg, is continuous on the space of local tensors. If $Q^2=I$, this map
    is an involution.
\end{theorem}

\begin{proof}\leanok
    Each tensor entry is a finite linear combination of the original entries.
    Applying two physical left actions gives
    $P(Q\mathcal U)=(PQ)\mathcal U$; hence $Q^2=I$ gives the involution.
\end{proof}

\begin{theorem}[Canonical form under a physical output unitary]
    \label{thm:mpu_canonical_form_physical_left_action}
    \lean{MPOTensor.isMPUCanonicalForm_ketLeftMul}
    \leanok
    \uses{def:mpu_canonical_form, def:mpdo_vertical_one_site_actions}
    Let $\mathcal U$ be a canonical-form MPO tensor, and let
    $Q\in\MN{d}$ be unitary. Then $Q\mathcal U$, with $Q$ acting on the
    physical output, is again in canonical form. Its canonical blocks
    have the same dimensions and weights, and the same ambient
    coisometry.
\end{theorem}

\begin{proof}\leanok
    \uses{thm:transfer_map_rotate_physical,
        thm:irreducible_tensor_rotate_physical,
        thm:rotate_physical_blocks}
    On the doubled physical alphabet, the action is the unitary rotation
    $Q\otimes I$. Apply this rotation to every canonical block. It
    preserves irreducibility and the transfer map, hence its spectral
    normalization, and distributes over the canonical block decomposition.
\end{proof}

\begin{theorem}[Transport of symmetry-preserving paths]
    \label{thm:mpu_symmetry_path_transport}
    \lean{MPOTensor.joinedIn_isMPU_isInvariantUnderSymmetry_iff_ketLeftMul,
        MPOTensor.strictlyEquivalentUnderSymmetry_iff_ketLeftMul}
    \leanok
    \uses{def:strictly-equivalent-symmetry,
        def:mpdo_vertical_one_site_actions}
    Fix $d$, two finite-chain operator symmetries $\mathcal S,\mathcal T$,
    and a unitary $Q\in\MN{d}$ with $Q^2=I$. Suppose that, for every local
    tensor $\mathcal W$ with physical dimension $d$ and any virtual
    dimension, $\mathcal W$ is $\mathcal S$-invariant if and only if
    $Q\mathcal W$ is $\mathcal T$-invariant. Then, for tensors
    $\mathcal U,\mathcal V$ of a common virtual dimension, a continuous MPU
    path of that virtual dimension from $\mathcal U$ to $\mathcal V$ with
    $\mathcal S$-invariance exists if and only if such a path from
    $Q\mathcal U$ to $Q\mathcal V$ with $\mathcal T$-invariance exists.
    For tensors $\mathcal U,\mathcal V$ of possibly different virtual
    dimensions, strict equivalence under $\mathcal S$ is equivalent to
    strict equivalence of $Q\mathcal U,Q\mathcal V$ under $\mathcal T$;
    both relations compare the tensors after unused virtual directions are
    adjoined up to a common dimension \cite{gap:mpu_equivalence_fixed_bond}.
    Identity-ancilla stabilization is outside the statement
    \cite{gap:mpu_symmetry_ancilla_transport}.
\end{theorem}

\begin{proof}\leanok
    \uses{thm:mpu_continuous_involutive_left_action,
        thm:threeMPU2_chain_formulas,
        thm:mpu_canonical_form_physical_left_action}
    If $p\mapsto\mathcal W(p)$ is the first path, then
    $p\mapsto Q\mathcal W(p)$ is continuous and has the required
    endpoints and symmetry. Its periodic operators are
    $Q^{\otimes N}W^{(N)}(p)$ and are unitary. Applying the same involution
    proves the converse. For strict equivalence, apply this to the paths
    between the enlarged tensors; the action of $Q$ commutes with adjoining
    unused virtual directions. Canonical form is preserved at both endpoints by
    Theorem~\ref{thm:mpu_canonical_form_physical_left_action}.
\end{proof}

\begin{theorem}[Strict swap comparison]
    \label{thm:mpu_strict_swap_comparison}
    \lean{MPOTensor.strictlyEquivalentUnderSymmetry_shiftSwapDagger_iff,
        MPOTensor.strictlyEquivalentUnderSymmetry_shiftSwapTranspose_iff,
        MPOTensor.strictlyEquivalentUnderSymmetry_conjugation_ketLeftMul_swap_iff}
    \leanok
    \uses{def:strictly-equivalent-symmetry,
        def:mpu_swap_comparison_symmetries, threeMPU2}
    Let $\mathcal U,\mathcal V$ have two $d$-dimensional physical spins,
    and write $\widetilde{\mathcal U}=\mathbb S\mathcal U$ and
    $\widetilde{\mathcal V}=\mathbb S\mathcal V$. Strict equivalence of
    $\mathcal U,\mathcal V$ under swap-combined adjunction is equivalent
    to strict equivalence of
    $\widetilde{\mathcal U},\widetilde{\mathcal V}$ under adjunction.
    The analogous equivalence holds for swap-combined transposition and
    transposition. Strict equivalence under conjugation is preserved.
    Identity-ancilla stabilization is outside this statement
    \cite{gap:mpu_symmetry_ancilla_transport}.
\end{theorem}

\begin{proof}\leanok
    \uses{thm:mpu_one_sided_swap_invariance,
        thm:mpu_symmetry_path_transport,
        thm:threeMPU2_chain_formulas}
    Apply Theorem~\ref{thm:mpu_symmetry_path_transport} with
    $Q=\mathbb S$. The swap is a unitary involution, and the three
    required invariance equivalences are precisely those of
    Theorem~\ref{thm:mpu_one_sided_swap_invariance}.
\end{proof}

\begin{theorem}[Physical left actions commute with blocking]
    \label{thm:mpu_physical_left_action_blocking}
    \lean{MPOTensor.blockTensor_ketLeftMul}
    \leanok
    \uses{def:blocked_tensor, def:mpdo_vertical_one_site_actions}
    For every local MPO tensor $\mathcal U$, one-site matrix $Q$, and
    $k\geq0$, output-leg multiplication commutes with blocking:
    \begin{align}
        (Q\mathcal U)^{[k]}
         & =Q^{\otimes k}\mathcal U^{[k]}.
        \label{eq:mpu_left_action_blocking}
    \end{align}
    The identity includes the empty block, with $Q^{\otimes0}=1$.
\end{theorem}

\begin{proof}\leanok
    Expand each output-leg action in the ordered product of $k$ local
    tensors. The coefficient of an intermediate physical configuration
    $(r_0,\ldots,r_{k-1})$ is
    $\prod_{n=0}^{k-1} Q_{i_n,r_n}$, which is the corresponding entry of
    $Q^{\otimes k}$. Summing over these configurations gives
    (\ref{eq:mpu_left_action_blocking}).
\end{proof}

\begin{theorem}[Symmetry path transport after positive blocking]
    \label{thm:mpu_positive_block_symmetry_path_transport}
    \lean{MPOTensor.exists_strictlyEquivalentUnderSymmetry_blockTensor_iff_ketLeftMul}
    \leanok
    \uses{def:mpu_blocked_finite_chain_symmetry,
        def:strictly-equivalent-symmetry,
        def:mpdo_vertical_one_site_actions}
    Fix $d$, finite-chain operator symmetries $\mathcal S,\mathcal T$,
    and a unitary involution $Q\in\MN{d}$. For every $k>0$, suppose that
    every local tensor $\mathcal W$ of physical dimension $d^k$ and any
    virtual dimension is invariant under the blocked symmetry
    $\mathcal S^{[k]}$ if and only if
    $Q^{\otimes k}\mathcal W$ is invariant under $\mathcal T^{[k]}$.
    Then the existence of a common positive blocking length making
    $\mathcal U^{[k]},\mathcal V^{[k]}$ strictly equivalent under
    $\mathcal S^{[k]}$ is equivalent to the existence of such a length
    making $(Q\mathcal U)^{[k]},(Q\mathcal V)^{[k]}$ strictly equivalent
    under $\mathcal T^{[k]}$. No identity ancillas are included in this
    comparison \cite{gap:mpu_symmetry_ancilla_transport}.
\end{theorem}

\begin{proof}\leanok
    \uses{thm:mpu_symmetry_path_transport,
        thm:mpu_physical_left_action_blocking}
    Fix a positive blocking length $k$. The matrix $Q^{\otimes k}$ is
    again a unitary involution, so apply
    Theorem~\ref{thm:mpu_symmetry_path_transport} in the
    blocked physical space. Identify the transformed endpoints using
    (\ref{eq:mpu_left_action_blocking}), and then take the existential
    quantifier over $k>0$.
\end{proof}

\begin{theorem}[Swap comparison after positive blocking]
    \label{thm:mpu_positive_block_swap_comparison}
    \lean{MPOTensor.exists_strictlyEquivalentUnderSymmetry_blockTensor_shiftSwapDagger_iff,
        MPOTensor.exists_strictlyEquivalentUnderSymmetry_blockTensor_shiftSwapTranspose_iff,
        MPOTensor.exists_strictlyEquivalentUnderSymmetry_blockTensor_conjugation_swap_iff}
    \leanok
    \uses{def:mpu_blocked_finite_chain_symmetry,
        def:mpu_swap_comparison_symmetries,
        def:strictly-equivalent-symmetry, threeMPU2}
    Let $\mathcal U,\mathcal V$ be tensors with two $d$-dimensional
    physical spins. Put
    $\widetilde{\mathcal U}=\mathbb S\mathcal U$ and
    $\widetilde{\mathcal V}=\mathbb S\mathcal V$. The existence of a
    positive common blocking length for strict equivalence of
    $\mathcal U^{[k]},\mathcal V^{[k]}$ under blocked swap-combined
    adjunction is equivalent to the existence of such a length for
    $\widetilde{\mathcal U}^{[k]},\widetilde{\mathcal V}^{[k]}$ under
    blocked adjunction. The analogous statements hold for transposition
    and conjugation. No ancilla enlargement is asserted
    \cite{gap:mpu_symmetry_ancilla_transport}.
\end{theorem}

\begin{proof}\leanok
    \uses{thm:mpu_positive_block_symmetry_path_transport,
        thm:mpu_pointwise_symmetry_invariance_comparison,
        thm:mpu_involution_combined_fixed_equation,
        thm:mpu_one_sided_swap_invariance,
        thm:threeMPU2_chain_formulas}
    Blocking identifies the sitewise swap with
    $B_N=(\mathbb S^{\otimes k})^{\otimes N}$. Reindexing the original
    $Nk$ physical sites preserves adjunction, transposition, and
    entrywise conjugation, and
    $B_N^2=I$, $B_N^\dagger=B_N^{\mathsf T}=B_N$, and
    $\overline{B_N}=B_N$. Hence on every blocked tensor the same
    one-sided involution gives the three symmetry-invariance
    correspondences. Apply
    Theorem~\ref{thm:mpu_positive_block_symmetry_path_transport} with
    $Q=\mathbb S$.
\end{proof}


\begin{lemma}[Transformation with swaps and symmetry equivalence]
    \label{lemma:sym-trafo-swap}
    \notready
    \discussion{5715}
    \uses{threeMPU, threeMPU2, def:equivalent-symmetry}
    Two MPUs $U_i^{(N)}$ and $U_j^{(N)}$ are in the same phase under time
    reversal combined with the local swap, transposition combined with the
    local swap, or conjugation if and only if the corresponding
    $\widetilde U_i^{(N)}$ and $\widetilde U_j^{(N)}$ are in the same phase
    under time reversal, transposition, or conjugation, respectively.
    % Source lines: paper_v2.tex 2045--2085.
\end{lemma}

\begin{lemma}[Admissible transformation with swaps and symmetry equivalence]
    \label{lem:mpu_admissible_swap_symmetry_equivalence}
    \notready
    \discussion{5715}
    % **Scope restriction (admissible phase comparison):** every blocked endpoint
    % and comparison path carries the datum of
    % Definition~\ref{def:mpu_admissible_equivalence_datum}. Documented in
    % docs/paper-gaps/mpu_canonical_form_full_support.tex.
    \uses{threeMPU, threeMPU2, def:mpu_admissible_strict_symmetry_equivalence, def:mpu_admissible_symmetry_equivalence, def:mpu_reduced_full_support_source_datum, def:mpu_admissible_equivalence_datum}
    Fix one of the three symmetry comparisons and let
    $p\mapsto\mathcal W(p)$ be the source-constructed interpolation used
    in that comparison. Suppose it carries an admissible symmetry-path datum of
    blocking length one. Thus the path $\mathcal W(p)$ and the four fixed
    comparison paths $\overline{\mathcal W(p)}$, $\mathcal W(p)^\sharp$,
    $\mathcal W(p)^{\mathsf T}$, and $\mathcal S[\mathcal W(p)]$ are simple and
    carry continuously varying reduced full-support source data. Suppose
    also that every one- and two-site endpoint tensor used in the local
    comparison carries its own datum. Then $U_i^{(N)}$ and $U_j^{(N)}$ are in
    the same admissible phase under time reversal combined with the local swap,
    transposition combined with the local swap, or conjugation if and only if
    the corresponding $\widetilde U_i^{(N)}$ and $\widetilde U_j^{(N)}$ are in
    the same admissible phase under time reversal, transposition, or
    conjugation, respectively.
    % Source lines: paper_v2.tex 2045--2085.
\end{lemma}

\begin{lemma}[Conjugation phases of $U_2$ and $U_3$]
    \label{prop:U2-U3-trivial-conjugation}
    \notready
    \discussion{5715}
    \uses{threeMPU, conjphs, prop:conjsym-strict-equiv}
    After blocking two sites, $U_2^{(N)}$ and $U_3^{(N)}$ are both in the
    trivial phase under conjugation.
    % Source lines: paper_v2.tex 2110--2134.
\end{lemma}

\begin{lemma}[Admissible conjugation phases of $U_2$ and $U_3$]
    \label{prop:mpu_admissible_u2_u3_conjugation}
    \notready
    \discussion{5715}
    % **Scope restriction (two-site blocked data):** both blocked tensors carry
    % reduced full-support source data. No preservation under blocking is
    % asserted. Documented in
    % docs/paper-gaps/mpu_canonical_form_full_support.tex.
    \uses{threeMPU, thm:mpu_admissible_conjugation_odd_rank, thm:mpu_admissible_conjugation_even_rank, def:mpu_reduced_full_support_source_datum, def:mpu_admissible_equivalence_datum, def:mpu_admissible_symmetry_equivalence}
    Suppose the two-site blocked tensors carry reduced full-support source data
    and the standard-form comparison paths carry admissible path data of
    blocking length one. Then the two-site blocked tensors of $U_2$ and
    $U_3$ are both in the trivial admissible strict phase under conjugation.
    Consequently $U_2$ and $U_3$ lie in the same admissible equivalence class.
    % Source lines: paper_v2.tex 2110--2134.
\end{lemma}

\begin{lemma}[Continuous unitary path to the factor swap]
    \label{thm:mpu-continuous-unitary-swap-path}
    \lean{Matrix.swapSymmetricProjection,
        Matrix.swapAntisymmetricProjection,
        Matrix.swapPath,
        Matrix.exists_continuous_unitary_swapPath}
    \leanok
    Let $\mathbb S$ denote the swap operator on
    $\mathbb C^d\otimes\mathbb C^d$. There is a continuous map
    $S\colon\mathbb R\to \mathrm{U}(d^2)$ such that
    \begin{align}
        \bigl(S(0),S(1)\bigr) & = \bigl(\Id,\mathbb S\bigr).
    \end{align}
    One concrete choice is
    \begin{align}
        S(x) & = P_{\mathrm{sym}} + e^{\mathrm{i}\pi x}P_{\mathrm{anti}},
    \end{align}
    where
    \begin{align}
        P_{\mathrm{sym}} & = \frac{\Id+\mathbb S}{2},
    \end{align}
    and
    \begin{align}
        P_{\mathrm{anti}} & = \frac{\Id-\mathbb S}{2}.
    \end{align}
    % Source lines: paper_v2.tex 2148--2151 and 2229--2234.
\end{lemma}

\begin{proof}
    \leanok
    The swap is a self-adjoint involution:
    \begin{align}
        \mathbb S^2 & = \Id,
    \end{align}
    and
    \begin{align}
        \mathbb S^\dagger & = \mathbb S.
    \end{align}
    Hence the two projections satisfy
    \begin{align}
        \bigl(P_{\mathrm{sym}}^2,P_{\mathrm{anti}}^2\bigr)
         & = \bigl(P_{\mathrm{sym}},P_{\mathrm{anti}}\bigr),
    \end{align}
    \begin{align}
        \bigl(P_{\mathrm{sym}}P_{\mathrm{anti}},
        P_{\mathrm{anti}}P_{\mathrm{sym}}\bigr) & = (0,0),
    \end{align}
    \begin{align}
        \bigl(P_{\mathrm{sym}}^\dagger,P_{\mathrm{anti}}^\dagger\bigr)
         & = \bigl(P_{\mathrm{sym}},P_{\mathrm{anti}}\bigr),
    \end{align}
    and
    \begin{align}
        P_{\mathrm{sym}}+P_{\mathrm{anti}} & = \Id.
    \end{align}
    Since $\lvert e^{\mathrm{i}\pi x}\rvert=1$, these identities imply
    \begin{align}
        S(x)S(x)^\dagger
         & = P_{\mathrm{sym}}+P_{\mathrm{anti}}=\Id.
    \end{align}
    Continuity follows from continuity of the exponential, while $e^0=1$ and
    $e^{\mathrm{i}\pi}=-1$ give the endpoint identities.
\end{proof}

\begin{lemma}[Time-reversal and transposition phases of $U_2$ and $U_3$]
    \label{prop:U2-U3-trivial-tr-transpose}
    \notready
    \discussion{5715}
    \uses{threeMPU2, lemma:sym-trafo-swap, thm:mpu-continuous-unitary-swap-path, def:strictly-equivalent-symmetry}
    The MPUs $\widetilde U_2^{(N)}$ and $\widetilde U_3^{(N)}$ are in the same
    phase under both time reversal and transposition.
    % Source lines: paper_v2.tex 2136--2155.
\end{lemma}

\begin{lemma}[Admissible time-reversal and transposition phases of $U_2$ and $U_3$]
    \label{prop:mpu_admissible_u2_u3_time_reversal_transposition}
    \notready
    \discussion{5715}
    % **Scope restriction (admissible symmetry path):** the comparison path
    % carries the datum of Definition~\ref{def:mpu_admissible_equivalence_datum}.
    % Documented in docs/paper-gaps/mpu_canonical_form_full_support.tex.
    \uses{threeMPU2, lem:mpu_admissible_swap_symmetry_equivalence, def:mpu_admissible_strict_symmetry_equivalence, def:mpu_reduced_full_support_source_datum, def:mpu_admissible_equivalence_datum}
    For each of time reversal and transposition, let
    $p\mapsto\mathcal W_{\mathcal S}(p)$ be the source-constructed
    interpolation from $\widetilde U_2$ to $\widetilde U_3$. Suppose it carries
    an admissible symmetry-path datum of blocking length one. Thus the path
    $\mathcal W_{\mathcal S}(p)$ and its fixed conjugate, physical-adjoint,
    transposed, and $\mathcal S$-image comparison paths are simple and carry
    continuously varying reduced full-support source data. Suppose also that
    the one- and two-site endpoint tensors used in the local comparison carry
    their own data. Then $\widetilde U_2^{(N)}$ and
    $\widetilde U_3^{(N)}$ are in the same admissible phase under both time
    reversal and transposition.
    % Source lines: paper_v2.tex 2136--2155.
\end{lemma}

\begin{lemma}[Conjugation phases of $U_1$ and $U_2$]
    \label{prop:eqiv-U1-U2-conjugate}
    \notready
    \discussion{5715}
    \uses{threeMPU2, lemma:sym-trafo-swap, conjphs, prop:conjsym-strict-equiv, cor:equiv-under-conjugation}
    The MPUs $U_1^{(N)}$ and $U_2^{(N)}$ lie in different strict conjugation
    phases exactly when $d=4k+2$ for a nonnegative integer $k$. Otherwise they
    lie in the same strict phase. All these phases are in Case I, so the two
    MPUs are equivalent after blocking and adjoining ancillas.
    % Source lines: paper_v2.tex 2158--2200.
\end{lemma}

\begin{lemma}[Admissible conjugation phases of $U_1$ and $U_2$]
    \label{prop:mpu_admissible_u1_u2_conjugation}
    \notready
    \discussion{5715}
    % **Scope restriction (admissible blockings and paths):** every standard form,
    % blocked endpoint, and comparison path carries the required source data.
    % Documented in docs/paper-gaps/mpu_canonical_form_full_support.tex.
    \uses{threeMPU2, lem:mpu_admissible_swap_symmetry_equivalence, thm:mpu_admissible_conjugation_odd_rank, thm:mpu_admissible_conjugation_even_rank, cor:mpu_admissible_equiv_under_conjugation, def:mpu_reduced_full_support_source_datum, def:mpu_admissible_equivalence_datum}
    For each endpoint tensor $\mathcal A$ among
    $U_1$, $U_2$, $\widetilde U_1$, and $\widetilde U_2$, suppose that
    $\mathcal A$, its conjugate comparison $\overline{\mathcal A}$, its
    two-site block $\mathcal A_2$, and the conjugate comparison
    $\overline{\mathcal A_2}$ are simple and carry their own reduced
    full-support source data. For every strict comparison used below, let
    $p\mapsto\mathcal W(p)$ be the
    source-constructed interpolation between the corresponding endpoint
    tensors. Suppose it carries an admissible conjugation-symmetry path datum of
    blocking length one: the path and its fixed conjugate, physical-adjoint,
    transposed, and conjugation-image comparison paths are simple and carry
    continuously varying source data. Require a length-one admissible path datum
    also for the blocked interpolation used in the final equivalence
    claim. Under these hypotheses,
    $U_1^{(N)}$ and $U_2^{(N)}$ are in different admissible strict phases under
    conjugation exactly when $k$ is a nonnegative integer and
    \begin{align}
        d & =4k+2.
    \end{align}
    Otherwise they are in the same admissible strict phase. All these standard
    forms lie in Case I, so $U_1^{(N)}$ and $U_2^{(N)}$ are admissibly equivalent
    after blocking and adding ancillas.
    % Source lines: paper_v2.tex 2158--2200.
\end{lemma}

\begin{definition}[Two-species ancilla configurations]
    \label{def:mpu_shift_ancilla_config}
    \lean{MPOTensor.ShiftAncillaConfig}
    \leanok
    \discussion{7470}
    For a periodic chain of length $N$, the configuration space consists of
    triples $((\sigma_1,\sigma_2),a)$ of two physical configurations and one
    $d$-dimensional ancilla configuration.
\end{definition}

\begin{definition}[Sitewise physical swap]
    \label{def:mpu_shift_ancilla_swap_12}
    \lean{MPOTensor.shiftAncillaSwap₁₂}
    \leanok
    \discussion{7470}
    \uses{def:mpu_shift_ancilla_config}
    The sitewise physical swap is
    \begin{align}
        S^{1,2}((\sigma_1,\sigma_2),a)
         & =((\sigma_2,\sigma_1),a).
    \end{align}
    % Source: paper_v2.tex 2217--2229 and Figure fig:TR-ancilla.
\end{definition}

\begin{definition}[Sitewise second-species--ancilla swap]
    \label{def:mpu_shift_ancilla_swap_2a}
    \lean{MPOTensor.shiftAncillaSwap₂a}
    \leanok
    \discussion{7470}
    \uses{def:mpu_shift_ancilla_config}
    The sitewise swap of the second species with the ancilla is
    \begin{align}
        S^{2,a}((\sigma_1,\sigma_2),a)
         & =((\sigma_1,a),\sigma_2).
    \end{align}
    % Source: paper_v2.tex 2217--2229 and Figure fig:TR-ancilla.
\end{definition}

\begin{definition}[Crossed first-species--ancilla swap]
    \label{def:mpu_shift_ancilla_swap_1a_next}
    \lean{MPOTensor.shiftAncillaSwap₁aNext,
        MPOTensor.shiftAncillaSwap₁aNext_apply}
    \leanok
    \discussion{7470}
    \uses{def:mpu_shift_ancilla_config, def:mpo_rotate_config}
    Let $R$ cyclically rotate configurations on the chain. The crossed swap is
    \begin{align}
        S^{1,a}_{+}((\sigma_1,\sigma_2),a)
         & =((Ra,\sigma_2),R^{-1}\sigma_1).
    \end{align}
    % Source: paper_v2.tex 2217--2229 and Figure fig:TR-ancilla.
\end{definition}

\begin{definition}[Explicit three-swap circuit]
    \label{def:mpu_shift_ancilla_three_swap}
    \lean{MPOTensor.shiftAncillaThreeSwap}
    \leanok
    \discussion{7470}
    \uses{def:mpu_shift_ancilla_swap_12,
        def:mpu_shift_ancilla_swap_2a,
        def:mpu_shift_ancilla_swap_1a_next}
    Let $C$ be the permutation obtained by applying $S^{1,a}_{+}$ first, then
    $S^{1,2}$, and finally $S^{2,a}$.
    % Source: paper_v2.tex 2217--2229 and Figure fig:TR-ancilla.
\end{definition}

\begin{definition}[Counter-shift permutation]
    \label{def:mpu_shift_ancilla_counter_shift}
    \lean{MPOTensor.shiftAncillaCounterShift}
    \leanok
    \discussion{7470}
    \uses{def:mpu_shift_ancilla_config, def:mpo_rotate_config}
    Define
    \begin{align}
        K((\sigma_1,\sigma_2),a)
         & =((R\sigma_2,R^{-1}\sigma_1),a).
    \end{align}
    % Source: paper_v2.tex 2217--2229 and Figure fig:TR-ancilla.
\end{definition}

\begin{lemma}[Three-swap permutation identity]
    \label{thm:mpu_shift_ancilla_three_swap_endpoint}
    \lean{MPOTensor.shiftAncillaThreeSwap_conj_swap₁₂,
        MPOTensor.shiftAncillaThreeSwap_conj_swap₁₂_apply,
        MPOTensor.permMatrix_shiftAncillaThreeSwap_endpoint,
        MPOTensor.permMatrix_shiftAncillaThreeSwap_conj_swap₁₂}
    \leanok
    \discussion{7470}
    \uses{threeMPU2,
        def:mpu_shift_ancilla_swap_12,
        def:mpu_shift_ancilla_three_swap,
        def:mpu_shift_ancilla_counter_shift}
    The three-swap circuit satisfies
    \begin{align}
        C^{-1}S^{1,2}C & =K.
    \end{align}
    Writing $P_f$ for the permutation matrix of $f$, one also has
    \begin{align}
        P_C P_{S^{1,2}} P_{C^{-1}} & =P_K.
    \end{align}
    At the coordinate level, $K$ is the configuration map of
    $\widetilde U_3$, whereas $\widetilde U_2$ has configuration map
    $((\sigma_1,\sigma_2),a)\mapsto
        ((R^{-1}\sigma_2,R\sigma_1),a)$.
    In \cite[Figure~16]{Cirac2017MPU}, the printed source labels produce a discrepancy at the $\widetilde U_2$
    boundary on a nondegenerate periodic
    chain.
    % Source: paper_v2.tex 2217--2229 and Figure fig:TR-ancilla.
\end{lemma}

\begin{proof}\leanok
    \uses{def:mpu_shift_ancilla_swap_12,
        def:mpu_shift_ancilla_three_swap,
        def:mpu_shift_ancilla_counter_shift}
    Evaluating the three swaps gives
    \begin{align}
        C^{-1}S^{1,2}C((\sigma_1,\sigma_2),a)
         & =((R\sigma_2,R^{-1}\sigma_1),a).
    \end{align}
    Extensionality gives the permutation equality. The permutation-matrix
    identity follows by functoriality.
\end{proof}

\begin{lemma}[Ancilla equivalence of $U_1$ and $U_2$]
    \label{prop:U1-U2-equiv-ancillatrick}
    \notready
    \discussion{5715}
    \uses{threeMPU2,
        lemma:sym-trafo-swap,
        thm:mpu-continuous-unitary-swap-path,
        def:strictly-equivalent-symmetry}
    After adding one ancilla of dimension $d$ per site,
    $\widetilde U_1^{(N)}$ and $\widetilde U_2^{(N)}$ are strictly equivalent
    under both time reversal and transposition.
    % Source lines: paper_v2.tex 2203--2235.
\end{lemma}

\begin{lemma}[Admissible ancilla equivalence of $U_1$ and $U_2$]
    \label{prop:mpu_admissible_u1_u2_ancilla}
    \notready
    \discussion{5715}
    % **Scope restriction (ancilla endpoint and admissible path data):** the
    % ancilla-enlarged tensors carry source data and the comparison path is
    % admissible. No preservation under adjoining ancillas is asserted.
    % Documented in docs/paper-gaps/mpu_canonical_form_full_support.tex.
    \uses{threeMPU2,
        lem:mpu_admissible_swap_symmetry_equivalence,
        def:mpu_admissible_strict_symmetry_equivalence,
        def:mpu_admissible_symmetry_equivalence,
        def:mpu_reduced_full_support_source_datum,
        def:mpu_admissible_equivalence_datum}
    After adding one ancilla of dimension $d$ per site and performing the source
    blocking, denote the enlarged blocked endpoints by
    $\widehat{\mathcal U}_1$ and $\widehat{\mathcal U}_2$. For either time
    reversal or transposition, let $p\mapsto\mathcal W_{\mathcal S}(p)$ be the
    source-constructed interpolation from $\widehat{\mathcal U}_1$ to
    $\widehat{\mathcal U}_2$. Suppose this interpolation carries an admissible
    symmetry-path datum of blocking length one. Thus the path
    $\mathcal W_{\mathcal S}(p)$ and its fixed conjugate, physical-adjoint,
    transposed, and $\mathcal S$-image comparison paths are simple and carry
    continuously varying reduced full-support source data. Suppose also that
    the one- and two-site enlarged endpoint tensors used in the local comparison
    carry their own data. Then $\widehat{\mathcal U}_1$ and
    $\widehat{\mathcal U}_2$ are admissibly strictly equivalent under both time
    reversal and transposition. No equivalence of the original ancilla-free
    tensors is asserted, because the equal ancilla dimensions used here need not
    satisfy the coprimality condition in
    Definition~\ref{def:mpu_admissible_symmetry_equivalence}.
    % Source lines: paper_v2.tex 2203--2235.
\end{lemma}

\subsection{Intrinsic blocked traces and the twisted transfer}

\begin{definition}[Twisted transfer of the squared operator tensor]
    \label{def:mpu_time_reversal_twisted_transfer}
    \lean{MPOTensor.timeReversalTwistedTransfer}
    \leanok
    Let $W$ be an MPO tensor with physical dimension $d$ and bond dimension
    $D$. In product bond coordinates, define
    \begin{align}
        T_W=\frac{1}{d}\sum_{i,j=0}^{d-1}W^{ij}\otimes W^{ji}.
    \end{align}
    Equivalently, this is the normalized physical diagonal of the product
    tensor for $W^2$. It is not assumed positive or Hermitian. The algebraic
    convention $0^{-1}=0$ permits zero physical dimension as well.
\end{definition}

\begin{lemma}[Opposite-word expressions for blocked exchange traces]
    \label{lem:mpu_time_reversal_blocked_open_words}
    \lean{MPOTensor.timeReversalPeriodicTrace_eq_sum,
        MPOTensor.timeReversalPeriodicTrace_blockTensor,
        MPOTensor.timeReversalTwistedTransfer_pow_eq_sum_evalWord}
    \leanok
    \uses{def:mpu_time_reversal_twisted_transfer}
    The exchange-contracted trace of the two-site periodic operator is
    \begin{align}
        \tau(W)=\sum_{i,j=0}^{d-1}\tr(W^{ij}W^{ji}).
    \end{align}
    For $k\geq0$, write $W^{ww'}$ for the ordered product along
    physical words $w,w'$ of length $k$. Then
    \begin{align}
        \tau(W_k)
         & =\sum_{w,w'}\tr(W^{ww'}W^{w'w}), \\
        T_W^k
         & =d^{-k}\sum_{w,w'}
        W^{ww'}\otimes W^{w'w}.
    \end{align}
    The products use exactly $k$ physical sites, without a further blocking.
    These are algebraic identities for arbitrary MPO tensors, including the
    empty words when $k=0$.
\end{lemma}

\begin{lemma}[Intrinsic blocked trace as a twisted-transfer trace]
    \label{lem:mpu_intrinsic_blocked_twisted_trace}
    \lean{MPOTensor.swapMatrix_commute_timeReversalTwistedTransfer,
        MPOTensor.intrinsicTimeReversalSign_blockTensor_eq_twistedTransfer_trace}
    \leanok
    \uses{def:mpu_time_reversal_twisted_transfer,
        lem:mpu_time_reversal_blocked_open_words}
    Let $F_D$ exchange the two bond factors. Then $F_DT_W=T_WF_D$, and the
    normalized intrinsic exchange trace of each block is
    \begin{align}
        \sigma_{\mathrm{int}}(W_k)=\tr(F_DT_W^k),\qquad k\geq0.
    \end{align}
    No canonical form, symmetry or standard-form choice is needed for these
    identities. They give a finite-dimensional expression relevant to
    \cite[lines~1633--1733]{Cirac2017MPU}.
    % The blocking parity proposition and the identification of the virtual
    % sign from even-chain comparison with the intrinsic sign are separate.
\end{lemma}

\begin{lemma}[Commuting weights remove stabilized nilpotent transients]
    \label{lem:mpu_stabilized_commuting_weighted_trace}
    \lean{Matrix.StabilizedRankOneData.isNilpotent_pow_sub_power,
        Matrix.StabilizedRankOneData.trace_mul_pow_eq_of_commute}
    \leanok
    Let $E$ be a finite complex matrix and suppose that, for some $J>0$,
    every power $E^n$ with $n\geq J$ equals a normalized rank-one projector
    $P=E^J$. Then $E^k-P$ is nilpotent for every $k>0$. If $F$ commutes
    with $E$, then
    \[
        \tr(FE^k)=\tr(FP),\qquad k>0.
    \]
    Indeed, $E^k-P=E^k(\Id-P)$ and the complementary idempotent shows
    that its $J$th power vanishes. Since $F$ commutes with this difference,
    their product is nilpotent and has trace zero. The hypotheses here
    describe a stabilized matrix; no assertion that a particular twisted
    transfer has these properties is made. This is the linear-algebra step
    needed to remove transient terms in a subsequent use of
    Lemma~\ref{lem:mpu_intrinsic_blocked_twisted_trace}.
\end{lemma}

\subsection{Intrinsic canonical blocking parity}

\begin{lemma}[Actual signed comparison of the twisted transfer]
    \label{lem:mpu_canonical_twisted_transfer_comparison}
    \lean{MPOTensor.IsMPUCanonicalFormII.exists_twistedTransfer_signed_unitary_comparison}
    \leanok
    \uses{lem:mpu_even_family_virtual_sign_gauge,
        def:mpu_time_reversal_twisted_transfer}
    Let $U$ be a canonical-form-II MPU tensor and suppose that
    $U^{(N)}=U^{(N)\dagger}$ for every even $N>1$.
    There are a sign $\eta\in\{1,-1\}$ and a virtual unitary $z$ such that,
    in the actual product bond coordinates,
    \[
        T_U=\eta(z^\dagger\otimes\Id)E_U(z\otimes\Id),
    \]
    where $E_U$ is the normalized completely positive transfer matrix.
    Indeed, the actual even-family virtual comparison gives
    $(U^\sharp)^{ij}=\eta zU^{ij}z^\dagger$, hence
    $U^{ij}=\eta z^\dagger(U^\sharp)^{ij}z$.
    Substitution into the normalized diagonal of the product tensor gives
    the displayed equality. The physical-adjoint double layer determines the
    stated product order; no exchange of bond factors is inserted.
\end{lemma}

\begin{lemma}[Derived intrinsic phase formula]
    \label{lem:mpu_canonical_intrinsic_blocking_phase_formula}
    \lean{MPOTensor.IsMPUCanonicalFormII.exists_intrinsic_blocking_phase_formula}
    \leanok
    \uses{lem:mpu_canonical_twisted_transfer_comparison,
        lem:mpu_stabilized_commuting_weighted_trace,
        lem:mpu_intrinsic_blocked_twisted_trace}
    Under the same canonical and even-chain adjoint-symmetry assumptions,
    there are $\eta\in\{1,-1\}$ and $C\in\mathbb C$ such that
    \begin{align}
        \sigma_{\mathrm{int}}(U)   & =\eta C,              \\
        \sigma_{\mathrm{int}}(U_k) & =\eta^k C\qquad(k>0).
    \end{align}
    The normalized MPU transfer stabilizes at a positive power to a
    normalized rank-one projector. Conjugating the bond-exchange weight
    through the preceding comparison gives a weight commuting with this
    actual transfer. Its nilpotent transient therefore has zero weighted
    trace at every positive power. Neither the comparison nor stabilization
    is an additional premise. The constant $C$ is retained, and $\eta$ is not
    identified with the first intrinsic value.
\end{lemma}

\begin{lemma}[Odd and even intrinsic blocking parity for canonical MPUs]
    \label{lem:mpu_canonical_intrinsic_blocking_parity}
    \lean{MPOTensor.IsMPUCanonicalFormII.intrinsic_blocking_trace_eq_of_odd,
        MPOTensor.IsMPUCanonicalFormII.intrinsic_blocking_trace_eq_of_even}
    \leanok
    \uses{lem:mpu_canonical_intrinsic_blocking_phase_formula}
    Let $U$ be a canonical-form-II MPU tensor with
    $U^{(N)}=U^{(N)\dagger}$ at all even $N>1$. For every $k>0$,
    \[
        \sigma_{\mathrm{int}}(U_k)=
        \begin{cases}
            \sigma_{\mathrm{int}}(U),   & k\text{ odd},  \\
            \sigma_{\mathrm{int}}(U_2), & k\text{ even}.
        \end{cases}
    \]
    This follows from $\eta^2=1$ in the preceding formula, with the two
    initial values retained separately. It does not assert that those
    values are equal or that the virtual sign equals the first value.
    This is intrinsic parity under explicit canonical assumptions; it does
    not construct the particular blocked gates or their sequential
    unwinding in \cite[lines~1643--1733]{Cirac2017MPU}. The unrestricted
    source proposition remains separate.
\end{lemma}

\begin{lemma}[Intrinsic blocked traces depend only on the physical family]
    \label{lem:mpu_intrinsic_blocked_trace_representative}
    \lean{MPOTensor.intrinsicTimeReversalSign_eq_of_mpo_two_eq,
        MPOTensor.intrinsicTimeReversalSign_blockTensor_eq_of_mpo_eq}
    \leanok
    Let $U,V$ have the same physical dimension, with possibly different bond
    dimensions. Equality $U^{(2)}=V^{(2)}$ gives
    $\sigma_{\mathrm{int}}(U)=\sigma_{\mathrm{int}}(V)$.
    If $U^{(N)}=V^{(N)}$ for every $N>0$, then for every $k>0$,
    \[
        \sigma_{\mathrm{int}}(U_k)=\sigma_{\mathrm{int}}(V_k).
    \]
    Indeed, the actual two-site periodic operators of the $k$-site blocks
    are the same coordinate reindexing of the original length-$2k$ operators.
    Equality at length zero and an identification of bond spaces are unnecessary.
\end{lemma}

\begin{lemma}[Intrinsic time-reversal blocking parity for arbitrary MPUs]
    \label{lem:mpu_intrinsic_blocking_parity}
    \lean{MPOTensor.IsMPU.intrinsic_blocking_trace_parity}
    \leanok
    \uses{lem:mpu_intrinsic_blocked_trace_representative,
        lem:mpu_canonical_intrinsic_blocking_parity,
        thm:mpu_reduced_cfii_representative}
    Let $d,D>0$ and let $U$ be an MPU tensor satisfying
    $U^{(N)}=U^{(N)\dagger}$ at every even $N>1$. For each $k>0$,
    \[
        \sigma_{\mathrm{int}}(U_k)=
        \begin{cases}
            \sigma_{\mathrm{int}}(U),   & k\text{ odd},  \\
            \sigma_{\mathrm{int}}(U_2), & k\text{ even}.
        \end{cases}
    \]
    Choose a reduced canonical representative with the same positive-length
    physical family. That family equality preserves the actual adjoint
    symmetry and all the intrinsic blocked traces. Apply the canonical
    result and transfer the identities back. Thus no canonical form,
    full-support datum or simplicity assumption is required for the input.
    This proves the intrinsic numerical parity pattern. The particular
    blocked gates in \cite[lines~1657--1729]{Cirac2017MPU} and their
    sequential unwinding remain a separate construction.
\end{lemma}

\begin{lemma}[Two-letter contraction from raw source cuts]
    \label{lem:mpu_raw_cut_two_letter_contraction}
    \lean{MPOTensor.mul_apply_eq_sum_of_sourceCut_factorizations}
    \leanok
    \uses{def:mpu_source_matrix_one, def:mpu_source_matrix_two}
    Suppose the two source cuts of $U$ factor as $M_1=X_1Y_1$ and
    $M_2=X_2Y_2$, with arbitrary finite intermediate spaces. Put
    \[
        A_{(l,r),(j_1,j_2)}
        =\sum_\beta (Y_2)_{l,(j_1,\beta)}(Y_1)_{r,(\beta,j_2)}.
    \]
    Then
    \[
        (U^{i_1j_1}U^{i_2j_2})_{\alpha\gamma}
        =\sum_{l,r}(X_2)_{(\alpha,i_1),l}
        A_{(l,r),(j_1,j_2)}(X_1)_{(i_2,\gamma),r}.
    \]
    This is a finite contraction identity. It needs no normalization,
    minimal cut dimensions, inverse matrices or canonical presentation.
\end{lemma}

\begin{definition}[The printed two-block gates]
    \label{def:mpu_printed_two_block_gates}
    \lean{MPOTensor.twoBlockInputGate, MPOTensor.twoBlockOutputGate}
    \leanok
    Let $u,v$ be square gates on $\mathbb C^d\otimes\mathbb C^d$.
    Grouping the four physical coordinates in the order $(12,34)$, define
    \[
        u^{(2)}=v_{23}(u\otimes u),\qquad
        v^{(2)}=v_{23}.
    \]
    Identity spectators occupy the first and fourth coordinates of the
    second expression. These are the gates chosen in
    \cite[lines~1657--1674]{Cirac2017MPU}.
\end{definition}

\begin{lemma}[Actual standard form with the printed two-block gates]
    \label{lem:mpu_supplied_printed_two_block_standard_form}
    \lean{MPOTensor.TwoSiteStandardFormData.blockTwo}
    \leanok
    \uses{def:mpu_printed_two_block_gates,
        lem:mpu_raw_cut_two_letter_contraction,
        lem:mpu_supplied_standard_form_source_cuts,
        def:mpu_supplied_two_site_standard_form}
    Let $W$ have a supplied standard form with square gates $u,v$ and
    positive physical and bond dimensions. Its actual two-site block has
    a supplied standard form with exactly the preceding gates $u^{(2)},v^{(2)}$.
    The half factors are the explicit left factors of the two raw cuts of
    $W$, with the canonical enumeration of their pair indices.
    The raw cut identities give its open tensor contraction; the two
    existing cut-contraction identities give $u^{(2)}$ and $v^{(2)}$.
    Their unitarity follows from that of $u,v$. No canonical presentation,
    normalized source datum, minimality or new compact decomposition is
    assumed. This constructs the blocking by two; arbitrary blocking and
    sequential unwinding remain separate.
\end{lemma}

\subsection{Actual finite Heisenberg action on complete blocks}
\label{sec:mpu_finite_heisenberg}

\begin{lemma}[Finite conjugation with arbitrary complete-block tail]
    \label{thm:mpu_finite_heisenberg_complete_block}
    \lean{MPOTensor.twoSiteHeisenbergInputEquiv,
        MPOTensor.twoSiteHeisenbergOutputEquiv,
        MPOTensor.twoSiteHeisenbergMatrix,
        MPOTensor.twoSiteHeisenbergMatrix_intertwiner,
        MPOTensor.twoSiteHeisenbergMatrix_isUnitaryBetween,
        MPOTensor.twoSiteHeisenbergMatrix_conjugation}
    \leanok
    \uses{def:mpu_supplied_two_site_standard_form,
        lem:mpu_two_site_open_three_block_buffers}
    Let $W$ have a supplied two-site standard form with gates $u,v$.
    For every $n\ge0$, write the actual periodic operator $W^{(3+n)}$
    in complete input coordinates separating the middle input pair,
    and complete output coordinates separating the four neighbouring
    output half-sites. Denote this reindexed operator by $M_n$.
    Then, for every middle-pair observable $A$,
    \begin{align}
        M_n(I\otimes A)            & =(I\otimes B_{u,v}(A))M_n,\notag \\
        M_n(I\otimes A)M_n^\dagger & =I\otimes B_{u,v}(A).
        \label{eq:mpu_finite_heisenberg_complete_block}
    \end{align}
    The exterior identity includes the whole arbitrary tail. The local
    output is therefore independent of its length. This is a finite
    complete-block consequence of
    the standard-form and finite-conjugation formulas in
    \cite[lines~603--622 and~2300--2306]{Cirac2017MPU}, under the scope restriction and remaining
    assembly recorded in \cite{gap:mpu_qca_complete_block_heisenberg_scope}.
\end{lemma}
\begin{proof}\leanok
    Fixing all exterior indices gives the already derived three-block
    trace intertwiner, with the arbitrary tail retained as a virtual word
    factor. Summing the middle indices identifies the two complete finite
    matrix products. The coordinate maps are bijective, and the actual
    periodic MPO is unitary by its supplied standard form. Multiplication
    by its adjoint gives the conjugation formula. No periodic seam or
    exterior word is removed in this argument.
\end{proof}

\begin{lemma}[Actual original-site conjugation with arbitrary residuals]
    \label{thm:mpu_finite_heisenberg_original_residuals}
    \lean{MPOTensor.originalSiteHeisenbergInputEquiv,
        MPOTensor.originalSiteHeisenbergOutputEquiv,
        MPOTensor.originalSiteHeisenbergMatrix,
        MPOTensor.originalSiteHeisenbergMatrix_intertwiner,
        MPOTensor.originalSiteHeisenbergMatrix_isUnitaryBetween,
        MPOTensor.originalSiteHeisenbergMatrix_conjugation}
    \leanok
    \uses{thm:mpu_original_site_buffered_three_block_locality,
        thm:mpu_finite_heisenberg_complete_block,
        def:mpu_supplied_two_site_standard_form}
    Let $U$ be an MPU, and suppose that the two-site block of its length-$k$
    block has supplied standard-form gates $u,v$. For arbitrary original
    residual lengths $p,s\ge0$ with $N=p+6k+s>1$, complete coordinates on
    the actual original periodic operator give
    \begin{align}
        M_{p,s}(I\otimes A)M_{p,s}^\dagger
         & =I\otimes B_{u,v}(A).
        \label{eq:mpu_finite_heisenberg_original_residuals}
    \end{align}
    The middle observable is on two complete length-$k$ blocks. The
    exterior identity includes every original-site residual coordinate.
    In particular, the local output is independent of both residual
    lengths, and no divisibility condition on $N$ is imposed.
    This is the one-window finite-conjugation consequence of
    \cite[lines~2300--2306]{Cirac2017MPU}, with the remaining scope recorded
    in \cite{gap:mpu_qca_complete_block_heisenberg_scope}.
\end{lemma}
\begin{proof}\leanok
    The actual direct-block equivalence identifies the three complete
    middle blocks with their standard-form letters. The original residual
    products remain at the virtual ends of the established three-block
    trace identity. Complete input and output coordinate bijections lift
    it to an intertwiner for the actual original finite operator. The MPU
    property at length $N>1$ gives its unitarity; multiplication by the
    adjoint proves the formula. No residual or periodic seam is discarded.
\end{proof}


\begin{lemma}[A fixed finite-window output under arbitrary buffering]
    \label{thm:mpu_finite_heisenberg_fixed_window}
    \lean{MPOTensor.bufferedWindowConfigEquiv,
        MPOTensor.bufferedWindowMatrix,
        MPOTensor.originalSiteHeisenbergInputWindow,
        MPOTensor.originalSiteHeisenbergOutputWindow,
        MPOTensor.mpo_bufferedWindow_conjugation}
    \leanok
    \uses{thm:mpu_finite_heisenberg_original_residuals}
    Under the preceding hypotheses, let $I_k(A)$ and $O_k(A)$ be the
    input and output matrices on the fixed window of $6k$ original
    sites. For every $p,s\ge0$ with $p+6k+s>1$,
    \[
        U^{(p+6k+s)}
        (I_p\otimes I_k(A)\otimes I_s)
        U^{(p+6k+s)\dagger}
        = I_p\otimes O_k(A)\otimes I_s .
    \]
    Both window matrices are independent of $p$ and $s$. Tensor factors
    are written in the original order of the sites. This is the finite
    buffering identity underlying
    \cite[lines~2300--2306]{Cirac2017MPU}; finite-region assembly is
    given below, with the inverse and completion scope recorded in
    \cite{gap:mpu_qca_complete_block_heisenberg_scope}.
\end{lemma}
\begin{proof}\leanok
    Separate the prefix and suffix words from the complete central
    window. In these coordinates, the exterior Kronecker identity is
    the product of the identities on the original residual words and
    on the exterior coordinates inside the window. Reindexing the
    actual finite conjugation identity gives the claimed formula.
\end{proof}


\begin{lemma}[The four consecutive middle output blocks]
    \label{lem:mpu_finite_heisenberg_middle_output_window}
    \lean{MPOTensor.wordOfBlock_twoSiteDirectBlockEquiv_symm,
        MPOTensor.originalSiteHeisenbergOutputWindowEquiv,
        MPOTensor.originalSiteHeisenbergOutputCentreEquiv,
        MPOTensor.originalSiteHeisenbergOutputWindowEquiv_middle,
        MPOTensor.originalSiteHeisenbergOutputWindowEquiv_outside_eq_iff,
        MPOTensor.originalSiteHeisenbergOutputWindow_eq_chainWindowOperator,
        MPOTensor.originalSiteHeisenbergOutputWindow_mem_supportedOperators}
    \leanok
    Let $k>0$. For arbitrary matrices $u,v,A$ of the standard-form gate
    dimensions, write $B_{u,v}(A)$ for the actual four-half-block output.
    The six-half-block output window is exactly
    \[
        O_k(A)=I_k\otimes B_{u,v}(A)\otimes I_k,
    \]
    where the middle four half blocks are decoded to their $4k$
    original sites in their given order. Thus $O_k(A)$ is the actual
    placement of that decoded matrix on the interval $[k,5k)$ of the
    $6k$-site window and is supported there. No output-support or
    factor-presentation condition is assumed. This is the literal
    four-half-block geometry of
    \cite[lines~603--622 and~2300--2334]{Cirac2017MPU}.
\end{lemma}
\begin{proof}\leanok
    \uses{thm:mpu_finite_heisenberg_fixed_window}
    The actual two-site physical identification decodes to the
    concatenation of its two original block words. The output coordinate
    order is therefore
    $(e_1,c_1,c_2,c_3,c_4,e_2)$, with $e_1,e_2$ the two exterior
    half blocks. Reading the interval $[k,5k)$ gives exactly the four
    middle coordinates. Equality outside it is equivalent to equality
    of the two exterior coordinates. In matrix entries, the exterior
    Kronecker identity is precisely the corresponding equality indicator.
    The resulting formula is the ordinary consecutive-window inclusion,
    which has the stated support.
\end{proof}

\begin{lemma}[Canonical local-observable buffering for one window]
    \label{thm:mpu_finite_heisenberg_canonical_window}
    \lean{SpinChain.intervalCoordinates_localInclusion_eq_bufferedWindowMatrix,
        SpinChain.finiteChainLocalObservable_originalSiteHeisenbergWindow}
    \leanok
    \uses{thm:mpu_finite_heisenberg_fixed_window}
    Fix an integer interval $\Lambda=[a,a+6k)$ and the preceding
    standard-form data. Interpret $I_k(A)$ and $O_k(A)$ as matrices
    on $\Lambda$ in consecutive site coordinates. For every $p,s\ge0$
    with $p+6k+s>1$, the actual finite-chain conjugation on
    $[a-p,a+6k+s)$, applied to the canonical inclusion of $I_k(A)$,
    represents exactly the local observable $O_k(A)$ in the algebraic
    local algebra. Consequently its value is independent of both
    buffers, including the position and length of the ambient interval.
    This is the one-window instance of
    \cite[lines~2300--2306]{Cirac2017MPU}; finite-region assembly is
    given below, with the inverse and completion scope recorded in
    \cite{gap:mpu_qca_complete_block_heisenberg_scope}.
\end{lemma}
\begin{proof}\leanok
    Equality of the complementary coordinates is precisely equality
    of both exterior words, so canonical local inclusion is the same
    identity-tensor inclusion in consecutive coordinates. Apply the
    preceding actual finite conjugation formula, then pass through
    the canonical inclusion into the algebraic local algebra. The
    latter identifies all enlargements of the fixed output interval.
\end{proof}


\begin{lemma}[The full central interval algebra has a buffered local action]
    \label{thm:mpu_finite_heisenberg_full_central_interval}
    \lean{MPOTensor.originalSiteHeisenbergCentreEquiv,
        SpinChain.intervalCoordinates_centralInclusion_originalSiteHeisenbergInputWindow,
        SpinChain.exists_finiteChainLocalObservable_central_interval}
    \leanok
    \uses{thm:mpu_finite_heisenberg_canonical_window}
    Under the preceding hypotheses, every observable on
    $\lbrack a+2k,a+4k)$ has a fixed output observable on
    $\lbrack a,a+6k)$.
    For every original-site buffering $\lbrack a-p,a+6k+s)$ of length
    greater than one, actual finite MPU conjugation of the canonical
    input inclusion represents this fixed output in the algebraic
    local algebra. Thus the complete algebra on the central $2k$
    sites is covered, rather than a chosen family of coefficient
    slices. Finite-region generation is derived below; the inverse
    family and completion remain separate from this instance of
    \cite[lines~2300--2306]{Cirac2017MPU}; see
    \cite{gap:mpu_qca_complete_block_heisenberg_scope}.
\end{lemma}
\begin{proof}\leanok
    The direct-block decoding bijection identifies the middle input
    coordinates with all configurations on the central $2k$ sites.
    Reading the three consecutive blocks shows that the middle
    restriction has exactly these coordinates, whereas equality
    outside it is exactly equality of the other two block indices.
    The input-window matrix is therefore the canonical inclusion of
    the complete central algebra. Apply the preceding buffered action
    and compose the two canonical input inclusions.
\end{proof}


\begin{lemma}[Actual buffering compatibility on arbitrary finite support]
    \label{thm:mpu_finite_heisenberg_arbitrary_region_compatibility}
    \lean{SpinChain.finiteChainLocalObservable_eq_of_twoSiteStandardForm_buffers}
    \leanok
    \uses{thm:mpu_finite_heisenberg_full_central_interval,
        lem:mpu_singleton_generation_and_inverse_reflection}
    Let $k>0$ and suppose that the length-$k$ blocked tensor has the
    supplied two-site standard form above. Let $\Lambda\subset\mathbb Z$
    be finite, and let $[b_j,b_j+N_j)$, for $j=1,2$, contain $\Lambda$,
    with $N_j>1$. Suppose that for every $x\in\Lambda$,
    \[
        b_j\le x-2k,
        \qquad x+4k\le b_j+N_j .
    \]
    Then the two actual finite-chain MPU conjugations give the same
    homomorphism from $\mathcal A_\Lambda$ into the algebraic local
    algebra. In particular a symmetric buffer of radius $4k$ suffices.
    This is derived compatibility for the supplied standard form in
    \cite[lines~2300--2306]{Cirac2017MPU}. Construction of that form from
    a bare MPU, the inverse family, completion, and the source radius
    remain separate as recorded in
    \cite{gap:mpu_qca_complete_block_heisenberg_scope}.
\end{lemma}
\begin{proof}\leanok
    For $x\in\Lambda$, place its one-site algebra in the complete central
    interval $[x,x+2k)$. The preceding theorem gives one fixed output
    on $[x-2k,x+4k)$ for both ambient chains. The two homomorphisms
    therefore agree on every one-site algebra in $\Lambda$. These
    algebras generate the complete finite-region algebra: disjoint
    regions generate their union by matrix tensor expansion, so
    induction over the sites places the whole algebra in the
    equalizer. For the empty region, both homomorphisms preserve its
    scalar algebra.
\end{proof}


\begin{lemma}[Actual compatible norm-preserving forward local action]
    \label{thm:mpu_finite_heisenberg_forward_local_action}
    \lean{SpinChain.finiteChainLocalObservable_localInclusion,
        SpinChain.bufferedMPULocalAction,
        SpinChain.exists_finiteChainLocalObservable_bufferedMPULocalAction,
        SpinChain.bufferedMPULocalAction_compatible,
        SpinChain.bufferedMPULocalAction_eq_finiteChainLocalObservable,
        SpinChain.bufferedMPUAlgebraicAction,
        SpinChain.norm_bufferedMPULocalAction,
        SpinChain.norm_bufferedMPUAlgebraicAction}
    \leanok
    \uses{thm:mpu_finite_heisenberg_arbitrary_region_compatibility}
    Under the preceding supplied standard-form hypotheses, choosing one
    sufficiently buffered finite chain for each region gives an actual
    compatible family of forward local actions. It therefore induces
    a unital star-algebra homomorphism of the algebraic local algebra.
    For positive physical dimension this homomorphism preserves the
    physical operator norm. The forward family is built from the
    actual finite MPU operators; compatibility and norm preservation
    are conclusions. This gives the forward algebraic-local part of
    \cite[lines~2292--2306]{Cirac2017MPU}; the inverse and completed
    automorphism remain separate as in
    \cite{gap:mpu_qca_complete_block_heisenberg_scope}.
\end{lemma}
\begin{proof}\leanok
    The largest absolute site coordinate of a finite support is finite.
    Adding $4k+1$ on both sides gives a canonical chain of length greater
    than one with all required buffers. For nested supports, the two
    chosen chain actions agree by the preceding compatibility theorem
    and composition of identity-tensor inclusions. The universal
    property of the algebraic local algebra therefore gives the forward
    homomorphism. Canonical local inclusion, actual unitary conjugation,
    and the inclusion into the algebraic local algebra all preserve the
    physical operator norm.
\end{proof}

\begin{lemma}[Completed actual isometric forward endomorphism]
    \label{thm:mpu_finite_heisenberg_completed_forward_action}
    \lean{SpinChain.bufferedMPUAlgebraicAction_isometry,
        SpinChain.bufferedMPUQuasiLocalAction,
        SpinChain.bufferedMPUQuasiLocalAction_apply,
        SpinChain.bufferedMPUQuasiLocalAction_algebraicToQuasiLocal,
        SpinChain.bufferedMPUQuasiLocalAction_quasiLocalObservable,
        SpinChain.bufferedMPUQuasiLocalAction_finiteChainConjugation,
        SpinChain.bufferedMPUQuasiLocalAction_isometry,
        SpinChain.norm_bufferedMPUQuasiLocalAction}
    \leanok
    \uses{thm:mpu_finite_heisenberg_forward_local_action}
    For positive physical dimension, the preceding actual forward action
    extends to an isometric unital star-algebra endomorphism of the
    quasi-local algebra. Its value on each finite observable is the
    canonical completed image of the actual buffered finite-chain
    conjugation already constructed. The same formula holds for every
    sufficiently buffered ambient interval, without imposing divisibility
    on its length. No inverse or surjectivity is
    assumed. This supplies the completed forward map in
    \cite[lines~2296--2306]{Cirac2017MPU}; the inverse and full source
    automorphism statement remain separate as in
    \cite{gap:mpu_qca_complete_block_heisenberg_scope}.
\end{lemma}
\begin{proof}\leanok
    Norm preservation makes the actual algebraic-local homomorphism an
    isometry. The uniform completion extends its ring homomorphism;
    preservation of complex scalars and the involution follow from the
    corresponding identities on the dense algebraic local algebra.
    Completion preserves the isometry. Evaluation on a finite observable
    follows from the extension formula on the canonical dense image.
\end{proof}

\begin{lemma}[Finite one-site generation and inverse support reflection]
    \label{lem:mpu_singleton_generation_and_inverse_reflection}
    \lean{SpinChain.range_localObservable_empty_le,
        SpinChain.range_localObservable_le_of_singletons,
        SpinChain.commute_localObservable_of_singletons,
        SpinChain.mem_range_localInclusion_iff_commute_outside_singletons,
        SpinChain.inverse_mem_range_localInclusion_of_forward_commute,
        SpinChain.inverse_mem_range_localInclusion_of_forward_support}
    \leanok
    The one-site algebras generate every finite local algebra; the empty
    local algebra consists of scalars. An observable on a finite region
    is supported in a retained subregion exactly when it commutes with
    every one-site algebra outside that subregion. Consequently, for a
    finite star-automorphism $e$, if $B$ commutes with the forward images
    of all one-site algebras outside a retained region $K$, then
    $e^{-1}(B)$ is supported in $K$. Forward support disjoint from the
    support of $B$ is sufficient for this commutation condition. These
    are finite algebraic ingredients for the actual inverse-locality
    argument in \cite[lines~2292--2306]{Cirac2017MPU}; no inverse family
    or global locality assertion is assumed or concluded here.
\end{lemma}
\begin{proof}\leanok
    Induct over disjoint unions of sites using matrix tensor expansion.
    Commutation with each generating algebra therefore implies
    commutation with the whole exterior algebra. Its relative commutant
    is exactly the retained identity-tensor factor. Applying $e^{-1}$
    to the forward commutators and using this characterization gives
    the finite inverse support assertion.
\end{proof}

\subsection{Finite inverse support with boundary collars}

\begin{lemma}[Actual finite inverse support with boundary collars]
    \label{lem:mpu_finite_inverse_buffered_region}
    \lean{SpinChain.finiteChainConjugation_symm_apply,
        SpinChain.finiteChainConjugation_symm_mem_range_of_twoSiteStandardForm_buffers,
        SpinChain.finiteInverseBufferedRegion,
        SpinChain.finiteChainConjugation_symm_mem_range_bufferedRegion}
    \leanok
    \uses{lem:mpu_singleton_generation_and_inverse_reflection}
    Let $U$ be an MPU with supplied two-site standard form after a
    positive blocking length $k$. Take a finite chain $[a,a+N)$ with
    $N>1$ and a target interval $[b,b+L)$ inside it. For every observable
    $B$ on that target, the actual inverse conjugation $U_N^\dagger B U_N$
    is supported in the intersection of the chain with
    \begin{align}
        \lbrack b-4k,b+L+4k)\;\cup\;\lbrack a,a+2k)\;\cup\;\lbrack a+N-4k,a+N).
        \label{eq:mpu_finite_inverse_collars}
    \end{align}
    The finite-chain boundary collars are retained in this statement.
\end{lemma}
\begin{proof}\leanok
    For any site outside (\ref{eq:mpu_finite_inverse_collars}), its
    forward standard-form window $[x-2k,x+4k)$ lies inside the chain
    and is disjoint from the target. The actual forward window identity
    gives the required commutator with $B$. Apply
    Lemma~\ref{lem:mpu_singleton_generation_and_inverse_reflection} to the actual
    finite unitary conjugation. Its inverse is precisely conjugation
    by the adjoint finite MPU.
\end{proof}

\begin{lemma}[Cyclic relocation of the actual finite conjugations]
    \label{lem:mpu_finite_conjugation_cyclic_relocation}
    \lean{MPOTensor.mpo_conjugation_submatrix_rotateConfig,
        MPOTensor.mpo_adjointConjugation_submatrix_rotateConfig,
        MPOTensor.mpo_conjugation_mem_supportedOperators_rotate,
        MPOTensor.mpo_adjointConjugation_mem_supportedOperators_rotate}
    \leanok
    Let $R$ simultaneously rotate the bra and ket configurations of a
    finite periodic chain. Both actual conjugations satisfy
    $R(U_N A U_N^\dagger)=U_N R(A)U_N^\dagger$ and
    $R(U_N^\dagger A U_N)=U_N^\dagger R(A)U_N$.
    Consequently any finite support bound for either conjugation
    rotates with the input through the periodic seam. Unitarity is
    not required for these covariance identities.
\end{lemma}
\begin{proof}\leanok
    The closed virtual trace gives $R(U_N)=U_N$. Simultaneous row and
    column relabeling preserves products and adjoints, yielding the
    two identities. A site permutation sends the finite supported
    operator algebra on $S$ to that on the rotated set $R(S)$.
\end{proof}

\begin{lemma}[Actual central-window action in the finite support algebras]
    \label{lem:mpu_finite_window_operator_support}
    \lean{MPSTensor.chainWindowOperator_eq_embedOp,
        MPSTensor.exists_chainWindowOperator_eq_of_mem_supportedOperators,
        MPOTensor.bufferedWindowMatrix_eq_chainWindowOperator,
        MPOTensor.bufferedWindowMatrix_mem_supportedOperators,
        MPOTensor.mpo_bufferedWindow_conjugation_mem_supportedOperators,
        MPOTensor.originalSiteHeisenbergInputWindow_eq_chainWindowOperator,
        MPOTensor.mpo_chainWindow_conjugation_mem_supportedOperators,
        MPOTensor.mpo_conjugation_mem_supportedOperators_central}
    \leanok
    \uses{thm:mpu_finite_heisenberg_fixed_window,
        thm:mpu_finite_heisenberg_full_central_interval}
    Under the supplied standard-form hypotheses, take a chain of length
    $N=p+6k+s$ with $k>0$. Every observable $C$ on its complete central
    interval $[p+2k,p+4k)$ has actual forward conjugate supported in
    $[p,p+6k)$. All exterior sites carry the identity, and $C$ is an
    arbitrary matrix on the complete $2k$-site space.
\end{lemma}
\begin{proof}\leanok
    For an operator $X$ on $L$ consecutive sites, write $W_{M,a}(X)$
    for $X$ placed on the sites $a,\dots,a+L-1$ of an $M$-site chain,
    with the identity elsewhere. Numbering an identity-tensor inclusion
    gives precisely this window operator, and placements compose:
    \begin{align}
        W_{M,a}\bigl(W_{w,b}(X)\bigr) & =W_{M,a+b}(X).
        \label{eq:mpu_finite_window_composition}
    \end{align}
    Let $\Phi$ be the bijection between the standard-form input
    coordinates and the configurations of the central interval. The
    source input window of $A$ is $I_k(A)=W_{6k,2k}(\Phi(A))$, so
    placing it at offset $p$ gives, by
    (\ref{eq:mpu_finite_window_composition}),
    \begin{align}
        W_{N,p}\bigl(I_k(A)\bigr) & =W_{N,p+2k}\bigl(\Phi(A)\bigr).
    \end{align}
    Every observable on $[p+2k,p+4k)$ has the form $W_{N,p+2k}(C)$, and
    $C=\Phi(A)$ for $A=\Phi^{-1}(C)$. The fixed-window formula of
    Lemma~\ref{thm:mpu_finite_heisenberg_fixed_window} then gives
    $U_N W_{N,p}(I_k(A))U_N^\dagger=W_{N,p}(O_k(A))$, which acts only on
    the $6k$ sites of $[p,p+6k)$.
\end{proof}


\begin{lemma}[Uniform actual finite-ring singleton propagation]
    \label{lem:mpu_finite_cyclic_singleton_propagation}
    \lean{QuantumCircuit.neighbourhood_finRotate,
        MPOTensor.mpo_conjugation_singleton_propagation_of_one,
        MPOTensor.mpo_conjugation_mem_supportedOperators_singleton}
    \leanok
    \uses{lem:mpu_finite_window_operator_support,
        lem:mpu_finite_conjugation_cyclic_relocation}
    Let an MPU have a supplied two-site standard form after a positive
    blocking length $k$. On every periodic ring of length $N\geq 6k$,
    the actual forward conjugate of an observable supported at a single
    site $x$ is supported in the ring neighbourhood of $x$ of radius $4k$.
    The same radius applies at every site, including the periodic seam.
\end{lemma}
\begin{proof}\leanok
    Write $\mathcal N_r(S)$ for the ring neighbourhood of radius $r$ of
    a set of sites $S$. First place $x=2k$ in the complete central
    interval of a $6k$-site window. The actual finite-window bound gives
    output support in $[0,6k)\subseteq\mathcal N_{4k}(\{2k\})$. For
    the unit rotation $R$, Lemma~\ref{lem:mpu_finite_conjugation_cyclic_relocation}
    and the rotation of neighbourhoods give
    \begin{align}
        R\bigl(U_NAU_N^\dagger\bigr)        & =U_NR(A)U_N^\dagger,       &
        R\bigl(\mathcal N_{4k}(\{x\})\bigr) & =\mathcal N_{4k}(\{x+1\}).
    \end{align}
    Thus the bound at $x$ gives the bound at $x+1$, with $x+1$ read
    modulo $N$. Finite cyclic induction from $x=2k$ gives the assertion
    at every site.
\end{proof}

\begin{lemma}[Actual translation covariance of the completed forward action]
    \label{lem:mpu_finite_heisenberg_translation_covariance}
    \lean{SpinChain.algebraicLocalTranslation_intervalCoordinates,
        SpinChain.quasiLocalTranslation_intervalCoordinates,
        SpinChain.intervalCoordinates_localInclusion_localTranslation,
        SpinChain.intervalCoordinates_finiteChainMPOMatrix,
        SpinChain.intervalCoordinates_finiteChainConjugation,
        SpinChain.algebraicLocalTranslation_finiteChainConjugation,
        SpinChain.quasiLocalTranslation_finiteChainConjugation,
        SpinChain.bufferedMPUQuasiLocalAction_translation_quasiLocalObservable,
        SpinChain.bufferedMPUQuasiLocalAction_translation}
    \leanok
    \uses{thm:mpu_finite_heisenberg_completed_forward_action}
    Translation of an interval observable shifts its integer anchor
    without changing its consecutive-coordinate matrix. The same is
    true of actual finite MPO conjugation. Consequently the completed
    forward action $F$ under the preceding supplied standard-form
    hypotheses commutes with every integer lattice translation:
    \[
        F(\tau_t(X))=\tau_t(F(X)),\qquad t\in\mathbb Z.
    \]
    Covariance is derived from the actual homogeneous finite operators;
    no covariance or inverse hypothesis is supplied. This is the
    forward translation statement in
    \cite[lines~2292--2306]{Cirac2017MPU}, with the source restriction
    recorded in \cite{gap:mpu_qca_complete_block_heisenberg_scope}.
\end{lemma}
\begin{proof}\leanok
    Canonical translation and consecutive coordinates give the same
    site numbering after shifting the interval anchor. Their local
    representatives agree after canonical inclusion. The periodic
    MPO matrix in these coordinates is independent of the anchor,
    so actual finite conjugation translates with its input. The
    translated interval satisfies the same buffer inequalities.
    Apply the already derived finite-chain evaluation of $F$ on both
    intervals. The resulting equality on finite local observables
    extends to the completion by continuity and density.
\end{proof}


\begin{lemma}[Actual forward propagation in the supplied block radius]
    \label{lem:mpu_finite_heisenberg_forward_propagation}
    \lean{SpinChain.bufferedMPUAlgebraicAction_supportedIn,
        SpinChain.bufferedMPUQuasiLocalAction_supportedIn}
    \leanok
    \uses{thm:mpu_finite_heisenberg_completed_forward_action,
        lem:mpu_singleton_generation_and_inverse_reflection}
    Under the preceding supplied positive block and standard-form
    hypotheses, the actual completed forward endomorphism $F$ sends
    every observable supported in $\Lambda$ to an observable supported
    in
    \[
        \Lambda+([-4k,4k]\cap\mathbb Z).
    \]
    The same support statement holds before completion in the algebraic
    local algebra. No output support or inverse property is supplied
    as a hypothesis. This is the forward propagation ingredient of
    \cite[lines~2298--2306]{Cirac2017MPU}, with the radius expressed
    in the supplied block length. Its bond-dimension bound and inverse
    assembly remain separate as in
    \cite{gap:mpu_qca_complete_block_heisenberg_scope}.
\end{lemma}
\begin{proof}\leanok
    For a one-site observable at $x$, choose the actual chain
    $[x-2k,x+4k)$ of length $6k>1$. The established finite-chain
    formula represents its output in precisely that finite-region
    algebra. This window is contained in the stated neighbourhood
    of $x$. The full one-site algebras generate each finite supported
    algebra; applying the forward homomorphism therefore places
    its whole image in the neighbourhood algebra. The completed
    evaluation formula gives the same support statement for every
    finite supported quasi-local observable.
\end{proof}


\begin{lemma}[Actual inverse propagation through the periodic seam]
    \label{lem:mpu_finite_cyclic_inverse_propagation}
    \lean{SpinChain.intervalCoordinates_localInclusion_eq_embedOp,
        SpinChain.intervalCoordinates_localInclusion_mem_supportedOperators,
        SpinChain.finiteRingNeighbourhood,
        SpinChain.finiteRingNeighbourhood_subset,
        SpinChain.finiteChainConjugation_symm_mem_range_ringNeighbourhood}
    \leanok
    \uses{lem:mpu_finite_cyclic_singleton_propagation,
        lem:mpu_singleton_generation_and_inverse_reflection}
    Under the supplied positive-block standard-form hypotheses, let
    $N\geq 6k$ and let $B$ be supported on an arbitrary finite target
    $J$ in the actual periodic chain of length $N$. Then the actual
    inverse conjugate $U_N^\dagger B U_N$ is supported in the ring
    neighbourhood of $J$ of radius $4k$. No boundary collar is needed.
\end{lemma}
\begin{proof}\leanok
    Consecutive site numbering sends a canonical local inclusion to
    its finite identity-tensor site embedding, with precisely the
    same support. Let $y\notin\mathcal N_{4k}(J)$ and let $C$ act at
    the single site $y$. By
    Lemma~\ref{lem:mpu_finite_cyclic_singleton_propagation},
    $U_NCU_N^\dagger$ is supported in $\mathcal N_{4k}(\{y\})$, which
    is disjoint from $J$. Hence
    \begin{align}
        \bigl[U_NCU_N^\dagger,B\bigr]=0
         & \quad\Longrightarrow\quad
        \bigl[C,U_N^\dagger BU_N\bigr]
        =U_N^\dagger\bigl[U_NCU_N^\dagger,B\bigr]U_N=0.
    \end{align}
    Thus $U_N^\dagger BU_N$ commutes with every one-site algebra outside
    $\mathcal N_{4k}(J)$. By
    Lemma~\ref{lem:mpu_singleton_generation_and_inverse_reflection}, it
    is supported in $\mathcal N_{4k}(J)$.
\end{proof}


\begin{lemma}[Actual completed automorphism from one supplied standard form]
    \label{lem:mpu_supplied_standard_form_actual_qca}
    \lean{SpinChain.finiteRingNeighbourhood_subset_regionSumset,
        SpinChain.exists_local_preimage_bufferedMPUQuasiLocalAction,
        SpinChain.bufferedMPUQuasiLocalAction_surjective,
        SpinChain.bufferedMPUQuasiLocalAut,
        SpinChain.bufferedMPUQuasiLocalAut_apply,
        SpinChain.bufferedMPUQuasiLocalAut_propagatesWithin,
        SpinChain.bufferedMPUQuasiLocalAut_isQCA}
    \leanok
    \uses{lem:mpu_finite_cyclic_inverse_propagation,
        lem:mpu_finite_heisenberg_forward_propagation,
        lem:mpu_finite_heisenberg_translation_covariance,
        thm:mpu_finite_heisenberg_forward_local_action,
        thm:mpu_finite_heisenberg_completed_forward_action}
    Under the supplied positive-block standard-form hypotheses, every
    observable on a finite target $J$ has a preimage under the actual
    completed forward action, supported in $J+[-4k,4k]$. Consequently
    that action is a star-automorphism of the quasi-local algebra,
    is translation covariant, and propagates by at most $4k$ sites.
    This is the supplied-standard-form specialization of the
    source construction of a quantum cellular automaton; it requires
    no inverse standard form.
\end{lemma}
\begin{proof}\leanok
    Choose one actual periodic chain buffering $J\cup\{0\}$ by six
    blocks on the left and twelve on the right. Its length is at
    least $6k$. The actual finite inverse of the target is supported
    in its ring neighbourhood of radius $4k$. The buffer prevents
    wraparound, so this neighbourhood lies in the ordinary integer
    neighbourhood $J+[-4k,4k]$, which itself has the buffers needed
    for actual forward evaluation. Finite unitarity makes that
    evaluation equal to the original target. Thus the range of the
    completed forward isometry contains every local observable.
    The range is closed by completeness and dense by the local
    observables, hence is the whole quasi-local algebra. The already
    derived forward covariance and propagation make it a quantum
    cellular automaton.
\end{proof}


\section{Finite support factors in the MPU--QCA index comparison}

\begin{lemma}[Least left support of the actual two-link output family]
    \label{lem:mpu_actual_two_link_left_support}
    \lean{Matrix.IsUnitaryBetween.conjStarAlgEquiv,
        Matrix.IsUnitaryBetween.conjStarAlgEquiv_apply,
        MPOTensor.twoLinkMiddleLift,
        MPOTensor.twoLinkOutputMap,
        MPOTensor.twoLinkOutputMap_apply,
        MPOTensor.twoLinkLeftOutputMap,
        MPOTensor.twoLinkRightOutputMap,
        MPOTensor.twoLinkOutputMap_inverse_conj_kronecker,
        MPOTensor.leftSupportAlgebra_twoLinkOutputMap,
        MPOTensor.twoLinkLeftSupportEquiv,
        MPOTensor.TwoSiteStandardFormData.left_div_phys_eq_phys_div_right}
    \notready
    Let $d,r,\ell$ be positive and let
    $u:\mathbb C^d\otimes\mathbb C^d\to\mathbb C^\ell\otimes\mathbb C^r$
    and $v:\mathbb C^r\otimes\mathbb C^\ell\to\mathbb C^d\otimes\mathbb C^d$
    be unitary. The actual output family
    \[
        A\longmapsto(v\otimes v)
        (I_r\otimes uAu^\dagger\otimes I_\ell)(v\otimes v)^\dagger
    \]
    is a star-homomorphism. The dimensions satisfy $\ell r=d^2$, so
    $\ell/d=d/r$. Its least left support algebra is
    $v(I_r\otimes M_\ell)v^\dagger$, and hence has matrix size $\ell$.
    This is the finite output-family step in
    \cite[Appendix, lines~2322--2334]{Cirac2017MPU}.
    The adjoint order follows the main-text direction of $v$; the
    Appendix correction and its distinct logarithmic factor-of-two
    correction are recorded in \cite{gap:mpu_qca_index_normalization}.
\end{lemma}
\begin{proof}
    Change the arbitrary input by the inverse of $u$. An intermediate
    product $A_\ell\otimes B_r$ then gives exactly
    $v(I_r\otimes A_\ell)v^\dagger\otimes
        v(B_r\otimes I_\ell)v^\dagger$. Matrix units span the intermediate
    algebra, so every left coefficient belongs to the first factor.
    Conversely choose $B_r=I_r$; a diagonal right coefficient recovers
    every element of the first factor. Leastness gives equality.
    Conjugation by $v$ and the faithful identity-tensor embedding
    identify it with $M_\ell$.
\end{proof}

\begin{lemma}[Coordinate transport of least finite left support]
    \label{lem:mpu_qca_left_support_coordinate_transport}
    \lean{Matrix.bipartiteSlice_reindex_prod,
        Matrix.leftSupportAlgebra_map_matrixReindexStarAlgEquiv}
    \notready
    Independent coordinate bijections on the two factors transport the
    least left support algebra by the first-factor bijection.
\end{lemma}
\begin{proof}
    A left coefficient in the new coordinates is exactly the
    first-factor coordinate transport of the old coefficient, with
    its two right indices relabelled. Apply the least-support
    characterization in both directions.
\end{proof}

\begin{lemma}[Intersection of finite quasi-local supports]
    \label{lem:mpu_qca_support_intersection}
    \lean{SpinChain.commute_map_localObservable_of_singletons,
        SpinChain.QuasiLocalSupportedIn.inter}
    \notready
    An observable supported on two finite regions $\Lambda$ and $\Gamma$
    is supported on $\Lambda\cap\Gamma$.
\end{lemma}
\begin{proof}
    Any site outside the intersection lies outside at least one
    of the two supports. The observable therefore commutes with
    its one-site algebra. Finite one-site generation gives commutation
    with every algebra disjoint from the intersection, and the
    outside-commutant characterization recovers exact support.
\end{proof}

\begin{lemma}[Actual complete-interval action and its four-block output]
    \label{lem:mpu_actual_complete_interval_four_output}
    \lean{SpinChain.bufferedMPUQuasiLocalAut_central_interval,
        SpinChain.bufferedMPUQuasiLocalAut_central_interval_four_output}
    \notready
    \uses{lem:mpu_supplied_standard_form_actual_qca,thm:quasi_local_bulk_inclusion}
    Under the supplied positive-block standard-form hypotheses, the
    actual completed action sends every matrix on the complete
    original-site interval $[a+2k,a+4k)$ to the two-link output matrix
    on $[a+k,a+5k)$, with the indicated block-coordinate identifications.
    In particular, its output support is this four-block interval.
\end{lemma}
\begin{proof}
    Choose the actual ambient chain $[a,a+8k)$ so that every input
    site has the required buffers. The finite-chain evaluation of
    the completed action and the actual standard-form identity give
    the fixed six-block output. Remove the two identity exterior
    blocks using its exact four-block placement and canonical
    interval inclusions.
\end{proof}


\begin{lemma}[Actual nearest-neighbour action after the supplied blocking]
    \label{lem:mpu_actual_blocked_nearest_neighbour}
    \lean{MPOTensor.originalSiteHeisenbergCentreEquiv_eq_blockedConfigEquiv,
        MPOTensor.originalSiteHeisenbergOutputBlockEquiv,
        SpinChain.blocked_bufferedMPUQuasiLocalAut_pair_interval,
        SpinChain.blocked_bufferedMPUQuasiLocalAut_pair_supportedIn,
        SpinChain.blocked_bufferedMPUQuasiLocalAut_singleton_supportedIn,
        SpinChain.blocked_bufferedMPUQuasiLocalAut_propagatesWithin}
    \notready
    \uses{lem:mpu_actual_complete_interval_four_output,
        lem:mpu_qca_support_intersection}
    Let $\omega$ be the actual supplied-standard-form automorphism and
    group its original sites in blocks of length $k>0$. The resulting
    automorphism has physical cell size $d_k=d^k$. Its actual action
    on every two-cell interval $[a,a+2)$ is the two-link output matrix
    on $[a-1,a+3)$. It has nearest-neighbour propagation on the blocked
    chain, without a supplied locality hypothesis.
\end{lemma}
\begin{proof}
    The actual blocking equivalence transports the complete two-block
    input and four-block output matrices to consecutive cell coordinates.
    A singleton belongs to either of its two adjacent input pairs.
    Intersecting their output regions gives exactly its three-cell
    neighbourhood. Finite one-site generation then gives this same
    nearest-neighbour propagation for every finite support.
\end{proof}

\begin{lemma}[The actual blocked MPU has GNVW index $\ell/d_k$]
    \label{lem:mpu_actual_blocked_gnvw_index}
    \lean{SpinChain.blocked_bufferedMPUQuasiLocalAut_evenSupportEquiv_zero,
        SpinChain.blocked_bufferedMPUQuasiLocalAut_supportAlgebraMatrixSize_zero,
        SpinChain.coe_gnvwIndex_blocked_bufferedMPUQuasiLocalAut,
        SpinChain.coe_gnvwIndex_blocked_bufferedMPUQuasiLocalAut_eq_phys_div_right}
    \notready
    \uses{def:qca_consecutive_pair_coordinates,lem:mpu_actual_blocked_nearest_neighbour,
        lem:mpu_actual_two_link_left_support,
        lem:mpu_qca_left_support_coordinate_transport,
        thm:qca_gnvw_index_dimtransfer}
    Under the supplied positive-block standard-form hypotheses, the
    actual blocked automorphism has canonical even support factor
    $\mathcal R_0$ on the cell pair $\{-1,0\}$ isomorphic to $M_\ell$.
    Its chosen matrix size at this cut is therefore $\ell$, and its
    positive rational GNVW index is
    \[
        \operatorname{ind}_{\mathrm{QCA}}
        =\frac{\ell}{d_k}=\frac{d_k}{r},\qquad d_k=d^k.
    \]
    This is the actual support-factor comparison after the source's
    grouping of sites in \cite[Appendix, lines~2313--2334]{Cirac2017MPU}.
    Independence of arbitrary regrouping is not asserted here. The
    Appendix normalization correction is recorded in
    \cite{gap:mpu_qca_index_normalization}.
\end{lemma}
\begin{proof}
    The exact finite action identifies the canonical ordered even-pair
    image with the actual two-link family in its left and right
    configuration coordinates. The canonical bipartite embedding is
    faithful, so this is an equality of finite image algebras. Transport
    the derived least left support through those independent coordinate
    bijections to obtain its full-matrix presentation of size $\ell$.
    Matrix-size uniqueness computes the chosen GNVW size and ratio.
    The actual supplied gate identity $\ell r=d_k^2$ gives the second
    ratio presentation.
\end{proof}


\subsection{Arbitrary finite buffers of a directly supplied tensor}
\label{sec:mpu_direct_supplied_window}

Let $W$ have physical size $d^2$ and supplied two-site standard-form gates
$u:\mathbb C^d\otimes\mathbb C^d\to\mathbb C^\ell\otimes\mathbb C^r$
and $v:\mathbb C^r\otimes\mathbb C^\ell\to\mathbb C^d\otimes\mathbb C^d$.
No preimage of $W$ under tensor blocking is assumed. This is the finite
standard-form calculation underlying
\cite[equation~StandardForm, lines~603--622 and~2300--2306]{Cirac2017MPU},
in the explicitly supplied-data scope of
\cite{gap:mpu_qca_complete_block_heisenberg_scope}.

\begin{definition}[Complete supplied-window coordinates]
    \label{def:mpu_supplied_window_coordinates}
    \lean{MPOTensor.twoSiteSuppliedInputEquiv,
        MPOTensor.twoSiteSuppliedOutputEquiv,
        MPOTensor.twoSiteSuppliedHeisenbergMatrix}
    \leanok
    For arbitrary prefix and suffix lengths $p,s\ge0$, write $N=p+3+s$.
    The full input coordinates separate the $p$ prefix cells, the first
    and third input cells of the central window, the middle input
    half-pair, and the $s$ suffix cells. The full output coordinates
    separate the same prefix and suffix, the two exterior half-sites,
    and the four output half-sites. Let $e_i,e_o$ be these bijections
    onto the full configuration set, and let
    \begin{align}
        M_{p,s} & =e_o^{-1}\operatorname{MPO}_W(N)e_i
    \end{align}
    denote the periodic matrix in these two coordinate systems.
\end{definition}

\begin{lemma}[Full finite intertwining and conjugation with arbitrary buffers]
    \label{lem:mpu_direct_supplied_window_conjugation}
    \lean{MPOTensor.twoSiteSuppliedHeisenbergMatrix_intertwiner,
        MPOTensor.twoSiteSuppliedHeisenbergMatrix_isUnitaryBetween,
        MPOTensor.twoSiteSuppliedHeisenbergMatrix_conjugation,
        MPOTensor.mpo_suppliedWindow_conjugation}
    \leanok
    \uses{def:mpu_supplied_window_coordinates,thm:mpu_finite_heisenberg_complete_block}
    Let $\Phi_{u,v}(A)$ be the actual four-half-site output observable.
    For every input half-pair matrix $A$ and arbitrary $p,s\ge0$,
    \begin{align}
        M_{p,s}(I_{\rm in}\otimes A)
         & =(I_{\rm out}\otimes\Phi_{u,v}(A))M_{p,s},\notag \\
        M_{p,s}(I_{\rm in}\otimes A)M_{p,s}^*
         & =I_{\rm out}\otimes\Phi_{u,v}(A).
    \end{align}
    The matrix $M_{p,s}$ is unitary between the two complete coordinate
    spaces. In the original configuration coordinates,
    \begin{align}
        \operatorname{MPO}_W(N)\,
        e_i(I_{\rm in}\otimes A)e_i^{-1}\,
        \operatorname{MPO}_W(N)^*
         & =e_o(I_{\rm out}\otimes\Phi_{u,v}(A))e_o^{-1}.
    \end{align}
    The actual periodic matrix, both buffers and the periodic seam remain
    present. No output-locality identity or additional blocking is a premise.
\end{lemma}
\begin{proof}
    \leanok
    Attach the prefix and suffix words to the existing open three-block
    intertwiner and take their actual virtual trace. The complete
    coordinate bijections lift the resulting coefficient identity to
    the full finite matrix product. The supplied standard form makes
    the actual periodic matrix unitary for $N\ge3$; cancelling its
    adjoint gives conjugation. Reindex both matrix products to obtain
    the original-coordinate formula.
\end{proof}


\subsection{The named completed action on a directly supplied tensor cell}
\label{sec:mpu_supplied_completed_local_action}

Let $W$ have physical size $d^2$ and supplied standard-form gates $u,v$.
Let $e:\{1,\ldots,b\}\to\{1,\ldots,d^2\}$ be a bijective relabelling,
and let $W'=e^*W$ be the actual reindexed tensor. The datum is explicitly
supplied, as in \cite{gap:mpu_qca_complete_block_heisenberg_scope}; no
preimage of $W$ under tensor blocking is introduced.

\begin{lemma}[Physical relabelling and complete window coordinates]
    \label{lem:mpu_supplied_relabel_coordinates}
    \lean{MPOTensor.IsMPU.reindexPhysical,
        MPOTensor.windowInputHeadEquiv,MPOTensor.windowOutputHeadEquiv,
        MPOTensor.twoSiteSuppliedCellEquiv,
        MPOTensor.twoSiteSuppliedInputWindowEquiv,
        MPOTensor.twoSiteSuppliedOutputWindowEquiv,
        MPOTensor.twoSiteSuppliedInputWindow,
        MPOTensor.twoSiteSuppliedOutputWindow}
    \notready
    \uses{def:mpu_supplied_window_coordinates}
    Physical relabelling preserves the actual periodic unitaries, since
    it simultaneously reindexes the row and column configurations.
    The singleton coordinate associated with a half-pair $c$ is the
    constant one-cell configuration $e^{-1}(\operatorname{pair}(c))$.
    The input head coordinates send $((j_0,j_2),c)$ to
    \begin{align}
        \bigl(e^{-1}(j_0),e^{-1}(\operatorname{pair}(c)),e^{-1}(j_2)\bigr).
    \end{align}
    The output coordinates send exterior halves $(a_L,a_R)$ and four
    output halves $(x_1,x_2,x_3,x_4)$ to
    \begin{align}
        \bigl(e^{-1}(\operatorname{pair}(a_L,x_1)),
        e^{-1}(\operatorname{pair}(x_2,x_3)),
        e^{-1}(\operatorname{pair}(x_4,a_R))\bigr).
    \end{align}
    Let $I_e(A)$ and $O_e(A)$ be the three-cell matrices obtained by
    adjoining identity exteriors to $A$ and $\Phi_{u,v}(A)$ in these
    complete coordinates.
\end{lemma}
\begin{proof}
    \notready
    Reindex each actual periodic matrix through the sitewise bijection;
    simultaneous reindexing preserves unitarity. The complete head
    equivalences already used by the finite-window calculation give the
    displayed coordinates, followed by sitewise physical relabelling.
\end{proof}

\begin{lemma}[Canonical buffered conjugation after physical relabelling]
    \label{lem:mpu_supplied_relabel_buffered_conjugation}
    \lean{MPOTensor.mpo_suppliedBufferedWindow_conjugation,
        MPOTensor.bufferedWindowMatrix_reindexPhysical,
        MPOTensor.mpo_reindexPhysical_suppliedBufferedWindow_conjugation,
        MPOTensor.twoSiteSuppliedInputWindow_eq_bufferedCell}
    \notready
    \uses{lem:mpu_supplied_relabel_coordinates,lem:mpu_direct_supplied_window_conjugation}
    Write $B_{p,s}(X)=I_p\otimes X\otimes I_s$ for actual canonical
    finite-window inclusion. Sitewise physical relabelling commutes with
    $B_{p,s}$. For every $p,s\ge0$, with $N=p+3+s$,
    \begin{align}
        \operatorname{MPO}_{W'}(N)\,B_{p,s}(I_e(A))\,
        \operatorname{MPO}_{W'}(N)^* & =B_{p,s}(O_e(A)).
    \end{align}
    Moreover $I_e(A)$ is exactly the singleton input matrix with one
    identity cell adjoined on each side.
\end{lemma}
\begin{proof}
    \notready
    The complete supplied-window coordinate formula equals canonical
    buffering by evaluating both matrices in the prefix/head/suffix
    configuration bijection. Apply the sitewise physical relabelling
    to the actual periodic conjugation identity. It commutes with
    concatenation of finite configurations and hence with identity
    buffering. The input head dictionary identifies its middle cell
    with the singleton coordinate, while the two exterior cells are
    identity factors.
\end{proof}

\begin{lemma}[Actual singleton output of the named automorphism]
    \label{lem:mpu_supplied_named_singleton_action}
    \lean{SpinChain.mpuQuasiLocalEquiv_supplied_singleton}
    \notready
    \uses{lem:mpu_supplied_relabel_buffered_conjugation,
        thm:mpu_actual_qca}
    Assume $b>0$. Let $\alpha_{W'}$ be the named actual bare-MPU
    automorphism and let $C_e(A)$ be the singleton-cell input matrix in
    the coordinates above. Then
    \begin{align}
        \alpha_{W'}\bigl(\iota_{[0,1)}(C_e(A))\bigr)
         & =\iota_{[-1,2)}(O_e(A)).
    \end{align}
    This is the actual supplied four-half-site output on cells $-1,0,1$.
    No compatible action, extra blocking scale or inverse-action datum
    is assumed.
\end{lemma}
\begin{proof}
    \notready
    The supplied standard form gives the MPU property of $W$, and
    physical reindexing gives that of $W'$. Choose exterior lengths
    $p=s=4D^4+1$. The resulting chain contains the singleton's full
    bare-MPU propagation buffer and the three-cell output window. The
    actual finite evaluation of the named automorphism therefore
    applies. Canonical local inclusions identify the singleton input
    with the central three-cell input window; the preceding periodic
    conjugation identity gives the stated output after faithful
    finite-region inclusion into the quasi-local algebra.
\end{proof}

